% =================================================
% 文件: 性健康入门指南.tex
% 说明: 面向初次接触性健康话题读者的科普入门读物
%       与同目录《性医学基础》(专业教材) 互补
% 版本: v1.0
% =================================================

\documentclass[12pt,UTF8,openany]{ctexbook}

\usepackage[top=2.6cm,bottom=2.6cm,left=2.8cm,right=2.8cm]{geometry}
\usepackage{amsmath}
\usepackage{amssymb}
\usepackage{array}
\usepackage{tabularx}
\usepackage{multirow}
\usepackage{booktabs}
\usepackage{enumitem}
\usepackage{xcolor}
\usepackage{graphicx}
\usepackage{tikz}
\usetikzlibrary{positioning,calc,arrows.meta,shapes.geometric,matrix}
\definecolor{myblue}{HTML}{185FA5}
\definecolor{myblueL}{HTML}{E6F1FB}
\definecolor{mygreen}{HTML}{3B6D11}
\definecolor{mygreenL}{HTML}{C0DD97}
\definecolor{myteal}{HTML}{1D9E75}
\definecolor{myamber}{HTML}{854F0B}
\definecolor{myamberL}{HTML}{FAEEDA}
\definecolor{myred}{HTML}{A32D2D}
\definecolor{myredL}{HTML}{FCEBEB}
\definecolor{mypurple}{HTML}{534AB7}
\definecolor{mypurpleL}{HTML}{EEEDFE}
\definecolor{mycoral}{HTML}{D85A30}
\definecolor{mycoralL}{HTML}{FAECE7}
\definecolor{mypink}{HTML}{D4537E}
\definecolor{mypinkL}{HTML}{F4C0D1}
\definecolor{mypinkM}{HTML}{ED93B1}
\definecolor{mygray}{HTML}{888780}
\usepackage[colorlinks,linkcolor=blue,citecolor=blue,urlcolor=blue]{hyperref}

\setCJKfamilyfont{hei}{SimHei}
\newcommand{\myhei}{\CJKfamily{hei}}

\ctexset{
  chapter/name={第,章},
  chapter/number={\arabic{chapter}},
  section/format=\Large\bfseries,
}

\newcolumntype{C}{>{\centering\arraybackslash}X}

\title{\myhei\zihao{0} 性健康从入门到实践}
\author{面向初次接触者的科学与安全手册}
\date{\today}

\begin{document}

\maketitle

\begin{center}
\itshape\zihao{4}
—— 从零开始，建立安全、自愿、知性的健康观念 ——
\end{center}

\vspace{1em}

\noindent\textbf{阅读提示：} 本册是一份性健康教育科普读物，面向零基础读者，
重点在于科学认知、自我保护和负责任的态度，而非任何具体技巧。
若你还未成年，请在信任的家长、校医或正规医疗机构指导下阅读与决策。
涉及个人健康的具体问题，请以专业医生的诊断为准。

\tableofcontents

\chapter{为什么要读这本书}

对性好奇、想提前弄明白，是非常正常且负责任的表现。真正危险的不是"懂"，
而是"似懂非懂就行动"。性健康是整体健康的一部分，它关系到身体的安全、
心理的成熟，以及人与人之间的尊重。

\section{一个核心判断}

本书反复强调一条主线，可以浓缩为六个字：

\begin{quote}
\centering \myhei\zihao{3}\textcolor{blue}{安全 · 自愿 · 知情}
\end{quote}

\begin{itemize}
  \item \textbf{安全}：采取科学措施，避免意外怀孕与性传播感染。
  \item \textbf{自愿}：任何行为都建立在双方清醒、自愿的前提之上。
  \item \textbf{知情}：了解足够的知识，在明白后果的前提下做决定。
\end{itemize}

\section{本书与专业教材的关系}

同一书库中另有一册《性医学基础》，系统讲解生殖解剖、性激素、性功能障碍等
专业内容，面向医学与相关专业读者。本书则面向完全没有基础的普通人，
从"我该了解什么、该准备什么"的角度切入，两者是互补关系。
当你读完本书想深入学习时，可以再阅读专业教材。

\chapter{五个方面的准备}

面对"初次接触性，我该准备什么"这个问题，可以把它拆成五个相互支撑的方面。
先建立整体框架，再逐一落实。

\begin{table}[htbp]
  \centering
  \caption{性健康准备的五个方面}
  \begin{tabularx}{\textwidth}{@{}l l X@{}}
    \toprule
    \textbf{方面} & \textbf{关键词} & \textbf{要解决的核心问题} \\
    \midrule
    认知准备   & 了解       & 懂基本的生理结构与生殖常识，不被谣言误导 \\
    安全防护   & 保护       & 学会正确使用安全套，懂得避孕与防病 \\
    沟通与同意 & 尊重       & 明确意愿与边界，学会和伴侣沟通 \\
    健康维护   & 预防       & 疫苗、筛查、卫生，把风险挡在前面 \\
    心态准备   & 从容       & 不攀比、不勉强，尊重彼此节奏 \\
    \bottomrule
  \end{tabularx}
\end{table}

以下各章按这五个方面依次展开。

\chapter{认知准备：先懂 basics}

\begin{center}
\begin{tikzpicture}[scale=0.85,
  organ/.style={fill=mypinkL,draw=mypink,thin},
  lead/.style={mygray,thin,dashed},
  lab/.style={font=\footnotesize}
]
  % 内生殖器（左）
  \draw[organ] (1,2.2) ellipse (0.35 and 0.5);
  \draw[organ] (4.5,2.2) ellipse (0.35 and 0.5);
  \draw[mypink,thick] (2.3,2.4) to[out=150,in=100] (1.4,2.3);
  \draw[mypink,thick] (3.7,2.4) to[out=30,in=80] (4.6,2.3);
  \fill[mypinkL,draw=mypink] (3,2.5) .. controls (2.6,2.2) and (2.6,1.7) .. (3,1.5) .. controls (3.4,1.7) and (3.4,2.2) .. (3,2.5);
  \fill[mypinkL,draw=mypink] (2.9,1.45) -- (2.85,0.5) -- (3.15,0.5) -- (3.1,1.45) -- cycle;
  \draw[lead] (1,2.6) -- (0.4,3) node[lab,left]{卵巢};
  \draw[lead] (2,2.5) -- (1.3,3.2) node[lab,above]{输卵管};
  \draw[lead] (3.4,2) -- (4.4,2.7) node[lab,right]{子宫};
  \draw[lead] (3.1,1) -- (4.4,0.8) node[lab,right]{阴道};
  % 外生殖器（右）
  \begin{scope}[shift={(7,0)}]
    \draw[mypink,thick] (0.5,2.6) to[out=-95,in=95] (0.5,0.6);
    \draw[mypink,thick] (2.5,2.6) to[out=-85,in=85] (2.5,0.6);
    \draw[mypinkM,thick] (0.9,2.5) to[out=-90,in=90] (0.9,0.8);
    \draw[mypinkM,thick] (2.1,2.5) to[out=-90,in=90] (2.1,0.8);
    \fill[mypinkM,draw=mypink] (1.5,2.7) circle (0.15);
    \fill[mypink] (1.5,2.1) circle (0.08);
    \draw[organ] (1.5,1.6) ellipse (0.3 and 0.2);
    \draw[lead] (1.5,2.85) -- (1.5,3.3) node[lab,above,mypink]{阴蒂};
    \draw[lead] (0.5,2) -- (-0.2,2.2) node[lab,left]{大阴唇};
    \draw[lead] (1.1,1.9) -- (0.2,1.9) node[lab,left]{小阴唇};
    \draw[lead] (1.6,2.1) -- (2.6,2.3) node[lab,right]{尿道口};
    \draw[lead] (1.8,1.6) -- (2.7,1.5) node[lab,right]{阴道口};
  \end{scope}
\end{tikzpicture}
\par\vspace{0.5em}\noindent\footnotesize 图：女性生殖系统示意（左：内生殖器；右：外生殖器）。科普示意，非精确解剖。
\end{center}

\section{了解身体与生命周期}

无论男女，了解自己的生殖系统结构与基本功能是第一步。建议通过正规医学
教材、医院科普栏或权威健康网站学习，理解以下内容：

\begin{itemize}
  \item 男性和女性外生殖器、内生殖器的名称与基本功能。
  \item 女性月经周期与排卵的大致规律，以及"安全期"并非绝对安全的原因。
  \item 精液、阴道分泌物等正常生理现象，以及哪些变化属于异常、需要就医。
\end{itemize}

\section{怀孕是怎么发生的}

怀孕需要精子与卵子结合（受精）并在子宫着床。理解这一点，就能理解各种
避孕方法的原理：有的是阻挡精子（屏障法），有的是抑制排卵（激素法），
有的是阻止着床。方法不同，但目标一致——在"不打算要孩子"时避免意外。

\section{常见误区}

\subsection*{误区一："体外排精"很安全}
失败率较高，因为射精前分泌物中可能已含有精子，且时点难以精确控制。

\subsection*{误区二："第一次不会怀孕"}
只要有排卵和 unprotected 的行为，就有怀孕的可能，与第几次无关。

\subsection*{误区三："看起来健康就没病"}
许多性传播感染在早期没有明显症状，无法靠外表判断，只能通过检测确知。

\chapter{安全防护：最实际的一章}

\begin{center}
\begin{tikzpicture}[
  st/.style={draw=myblue,fill=myblueL,rounded corners=2pt,minimum width=1.9cm,minimum height=1.1cm,align=center,font=\small}
]
\node[st] (s1) {1\\检查\\有效期/尺寸};
\node[st,right=0.5cm of s1] (s2) {2\\小心撕开\\(别咬)};
\node[st,right=0.5cm of s2] (s3) {3\\戴在\\勃起上};
\node[st,right=0.5cm of s3] (s4) {4\\全程使用\\不翻转};
\node[st,right=0.5cm of s4] (s5) {5\\射精后\\按住退出};
\draw[-Stealth] (s1) -- (s2);
\draw[-Stealth] (s2) -- (s3);
\draw[-Stealth] (s3) -- (s4);
\draw[-Stealth] (s4) -- (s5);
\end{tikzpicture}
\par\vspace{0.5em}
\begin{tikzpicture}[font=\small]
  \node[anchor=west,font=\normalsize\bfseries] at (0,1.6){紧急避孕时间窗（越早越好）};
  \fill[mygreenL] (0,0.9) rectangle (6,1.2);
  \node[anchor=west,mygreen] at (6.2,1.05){$\le$72h\ 左炔诺孕酮};
  \fill[myblueL] (0,0.3) rectangle (10,0.6);
  \node[anchor=west,myblue] at (10.2,0.45){$\le$120h\ 乌利司坦/铜IUD(最可靠)};
  \draw[mygray] (0,-0.2) -- (10,-0.2);
  \draw[mygray] (0,-0.3) -- (0,-0.1) node[below,mygray,font=\footnotesize] at (0,-0.4){0h};
  \draw[mygray] (6,-0.3) -- (6,-0.1) node[below,mygray,font=\footnotesize] at (6,-0.4){72h};
  \draw[mygray] (10,-0.3) -- (10,-0.1) node[below,mygray,font=\footnotesize] at (10,-0.4){120h};
\end{tikzpicture}
\end{center}

这是全书最具操作性的部分。请认真读完，并落实到位。

\section{安全套：一举两得的工具}

安全套（避孕套）是\textbf{目前唯一能同时避孕和降低性传播感染风险}的常用工具。
据美国 CDC 资料，在血清状态不一致（一方感染、一方未感染）的异性恋伴侣中，
坚持使用安全套者感染 HIV 的风险比不使用者的伴侣低约 71\%--80\%；
对衣原体、淋病、乙肝、滴虫等其他感染也有明确的降低风险作用。
世界卫生组织同样指出，安全套是唯一兼具避孕与预防包括 HIV 在内的性传播感染功能的方法。

\subsection*{关于有效率的真实数字}

\begin{table}[htbp]
  \centering
  \caption{男用乳胶安全套的第一年失败率（避孕）}
  \begin{tabularx}{\textwidth}{@{}l C C@{}}
    \toprule
    \textbf{使用方式} & \textbf{第一年失败率} & \textbf{说明} \\
    \midrule
    典型使用（日常真实情况） & 约 15\% & 常因遗漏、佩戴不当所致 \\
    坚持且正确使用           & 约 2\%  & 接近理想效果 \\
    \bottomrule
  \end{tabularx}
\end{table}

可见安全套本身可靠，问题几乎总是出在"没有每次、正确地使用"上。
CDC 也指出，安全套失效多源于使用不一致或不正确，而非破裂——
乳胶安全套在性交与退出过程中的破裂率约为每一百个中两个。

\subsection*{正确使用的要点}

\begin{enumerate}[leftmargin=2em,itemsep=0.3em]
  \item 每次行为都用\textbf{新的}安全套，检查包装完整与有效期。
  \item 在生殖器发生任何接触\textbf{之前}佩戴：阴茎勃起后，捏住储精囊排出空气，再展开至根部。
  \item 全程使用，结束后在阴茎仍勃起时按住套口边缘一同退出，避免精液外溢。
  \item 用后打结或用纸巾包好丢入垃圾桶，不可重复使用。
  \item 若中途破裂或滑脱，应立即停止并更换新的安全套。
  \item 润滑剂只用\textbf{水基或硅基}；凡士林、婴儿油、护肤乳等油基会削弱乳胶、导致破裂。
  \item 对乳胶过敏者可选聚氨酯等合成材质安全套。
\end{enumerate}

\section{其他避孕方法}

\begin{table}[htbp]
  \centering
  \caption{常见避孕方法概览（ efficacy 为典型使用下的大致水平）}
  \begin{tabularx}{\textwidth}{@{}l X C@{}}
    \toprule
    \textbf{方法} & \textbf{特点} & \textbf{能否防性病} \\
    \midrule
    男用安全套   & 易获取、即时可用；需每次正确使用 & 能 \\
    女用安全套   & 由使用者自主掌握；同样为屏障法   & 能 \\
    短效口服避孕药 & 需每日按时服用；适合固定伴侣   & 不能 \\
    宫内节育器   & 长效可逆，由医生放置           & 不能 \\
    皮下埋植     & 长效可逆，由医生放置           & 不能 \\
    绝育手术     & 永久性的，适用于确定不再生育者 & 不能 \\
    安全期 / 体外排精 & 可靠性低，不推荐单独依赖 & 不能 \\
    \bottomrule
  \end{tabularx}
\end{table}

选择哪种方法，应结合健康状况、伴侣情况与个人需求，并咨询医生。
注意：\textbf{除安全套外，其他方法基本不能预防性传播感染}，
必要时需与安全套联合使用（双重保护）。

\section{万一发生了无保护行为}

\begin{itemize}
  \item \textbf{紧急避孕}：在无保护行为后应尽快（越早越好，通常在 72 小时内，
        部分方案可延长至 120 小时）咨询医生或药师使用紧急避孕药，
        它不能作为常规避孕手段。
  \item \textbf{检测}：如有感染风险，应在医生建议的时间窗后做性传播感染筛查；
        很多感染早期无症状，检测是唯一确知方式。
  \item \textbf{就医}：出现疼痛、异常分泌物、出血、发热等症状，请及时到正规医院就诊。
\end{itemize}

\chapter{沟通与同意：不可逾越的底线}

\section{什么算"同意"}

同意必须同时满足几个条件：\textbf{清醒}（未受酒精或药物影响）、
\textbf{自愿}（没有被强迫、欺骗或施压）、\textbf{具体}（针对当下这件事）、
\textbf{可撤回}（任何时候都可以改变主意）。沉默、犹豫、畏惧都不等于同意。

需要特别强调：性行为涉及法律年龄界限，各地规定不同。若你或对方未成年，
请务必先了解当地法律，并听取可信成年人的意见。

\section{如何与伴侣沟通}

 Opening the conversation 常常让人觉得尴尬，但这是成熟与负责的表现。
可以在轻松的场合，坦诚地谈：

\begin{itemize}
  \item 双方对避孕与防病的打算：是否使用安全套、是否做过检测。
  \item 彼此的意愿、顾虑与边界：什么可以做、什么不接受。
  \item 可以随时喊停的"安全词"或约定，让任何一方安心表达不适。
\end{itemize}

谈这些不会破坏气氛，反而会让双方更放松、更信任。真正尊重你的人，
会认真对待你的边界。

\chapter{健康维护：把风险挡在前面}

\section{疫苗}

\subsection*{HPV 疫苗}
人乳头瘤病毒（HPV）感染与宫颈癌等多种疾病相关，疫苗是有效的预防手段。
在我国，二价、四价、九价 HPV 疫苗的获批接种年龄已扩展至 9--45 岁
（具体以最新说明书和当地接种门诊为准）。CDC 建议 11--12 岁常规接种，
可从 9 岁开始；26 岁前未接种者建议补种；27--45 岁者可与医生共同决策。
总体原则是：\textbf{越早接种，保护效果越好}。

\subsection*{乙肝疫苗}
乙肝可通过性接触传播。CDC 建议所有未接种、未感染且有风险者接种乙肝疫苗。

\section{定期筛查}

\begin{itemize}
  \item 有固定性伴侣前，双方做一次性传播感染筛查，是负责任的做法。
  \item 女性按年龄与医生建议定期进行宫颈癌筛查（如 HPV 检测、TCT）。
  \item 很多感染无症状，"没感觉"不等于"没感染"。
\end{itemize}

\section{日常卫生}

保持基本清洁即可，不必过度。重点是：行为前后适当清洁、
不共用可能接触体液的个人物品、出现不适及时就医而不是自行乱用药。

\chapter{心态准备：从容比技巧更重要}

\begin{itemize}
  \item \textbf{不必比较}：没有"标准剧本"，每个人的节奏都不同。
  \item \textbf{允许笨拙}：紧张、不熟练都很正常，沟通与体谅远比表现重要。
  \item \textbf{可以不做}："不想"本身就是充分的理由，任何人都不应被勉强。
  \item \textbf{尊重彼此}：把对方当作有感受、有边界的人，而不是目标。
\end{itemize}

健康的性观念，本质是把自己和对方都当作值得尊重的人。

\chapter{去哪里获取可靠信息}

信息源决定了观念的质量。请优先选择权威渠道，远离以流量为目的的
"经验分享"和以表演为目的的成人内容——后者不是教材，会建立错误的期待。

\begin{table}[htbp]
  \centering
  \caption{推荐的可信信息渠道}
  \begin{tabularx}{\textwidth}{@{}l X@{}}
    \toprule
    \textbf{渠道} & \textbf{说明} \\
    \midrule
    世界卫生组织（WHO）   & 全球权威的性与生殖健康指南与数据 \\
    美国疾控中心（CDC）   & 性病预防、安全套使用、疫苗接种的循证资料 \\
    国家卫健委 / 各地疾控 & 国内政策、疫苗接种点与健康教育信息 \\
    正规医院妇科 / 泌尿外科 / 皮肤性病科 & 个体化的诊断、检测与咨询 \\
    权威医学工具书         & 如默沙东诊疗手册（MSD Manual）等 \\
    \bottomrule
  \end{tabularx}
\end{table}

\chapter{深入篇：高频疑问解答}
\label{chap:faq}

前面各章给出的是框架，这一章回答初学者最常问的几个具体问题。
顺序从"工具怎么用"到"方案怎么选"，再到"话怎么讲""事怎么办"。

\section{安全套怎么才算"用对了"}

多数失败并非产品质量问题，而是使用环节出错。对照下表自查：

\begin{table}[htbp]
  \centering
  \caption{安全套使用的常见错误与正确做法}
  \begin{tabularx}{\textwidth}{@{}X X@{}}
    \toprule
    \textbf{常见错误} & \textbf{正确做法} \\
    \midrule
    中途才戴、先接触后戴 & 生殖器接触之前就戴好 \\
    用牙齿/指甲撕开包装   & 沿锯齿边小心撕开，避免划破 \\
    没排空气               & 捏住顶端储精囊排出空气再展开 \\
    戴反了翻过来再用       & 戴反即换新，不翻面重复使用 \\
    用油基润滑（凡士林等） & 只用水基或硅基润滑剂 \\
    一个用到底、重复使用   & 每次每个行为都用新的 \\
    射精后不按住就退出     & 按住套口、趁勃起时退出 \\
    存放在高温/暴晒处      & 阴凉干燥处存放，看有效期 \\
    \bottomrule
  \end{tabularx}
\end{table}

\section{避孕方式怎么选}

没有"最好"的方法，只有"最适合当下情况"的方法。可以分三步判断。

\subsection*{第一步：看关系状态}

\begin{itemize}
  \item \textbf{新伴侣、非固定关系，或对方健康状况不明}：\textbf{必须用安全套}。
        这是唯一能降低性传播感染风险的工具，不能省。
  \item \textbf{稳定单一伴侣且双方已检测}：可在安全套之外，叠加更高效的方法，
        如短效口服避孕药、宫内节育器（IUD）、皮下埋植，形成"双保险"。
  \item \textbf{还没决定或只是偶尔}：安全套通常已足够，不必把问题复杂化。
\end{itemize}

\subsection*{第二步：看有效率需求}

若把"意外怀孕"的风险降到最低，长效可逆避孕（IUD、埋植）与绝育术的有效率
最高，短效药次之，安全套与行为类方法相对较低。但请记住：
\textbf{高效避孕 $\neq$ 防病}，只有安全套兼顾后者。

\subsection*{第三步：看健康与可及性}

激素类方法并非人人适合（如有某些血栓风险、偏病史等情况需医生评估）；
IUD、埋植需由医生放置；短效药需每日按时服用。请结合身体状况咨询医生，
不要自行长期服用来源不明的药物。

\section{怎么和伴侣开口谈检测和边界}

很多人觉得难以启齿，但这恰恰是成熟的表现。可以事先想好一句开场白，例如：

\begin{quote}
"我想认真聊一下我们的健康和安全——不是不信任你，
而是我们都该对自己负责。我们要不要一起去做个检测，
以及确认一下彼此都能接受的边界？"
\end{quote}

要点：\textbf{对事不对人}，强调"保护我们俩"而非"怀疑你"；
把话题落在具体行动（检测、安全套、可以喊停）上，而不是抽象审问。
对方若尊重你，会认真对待；若回避、施压或嘲讽，这本身就是重要信号。

\section{无保护行为后怎么办}

\begin{enumerate}[leftmargin=2em,itemsep=0.3em]
  \item \textbf{尽快求助}：紧急避孕越早越好，通常在 72 小时内，部分方案可延至 120 小时，
        请咨询医生或药师，不要当作常规避孕手段。
  \item \textbf{安排检测}：很多感染早期无症状，需在医生建议的时间窗后筛查；
        过早检测可能出现假阴性。
  \item \textbf{留意症状}：出现异常分泌物、疼痛、出血、发热等，及时就医。
  \item \textbf{下次改进}：把这次经历转化为经验，落实"每次、全程、正确"使用安全套。
\end{enumerate}

\section{什么时候必须去看医生}

出现以下任一情况，请尽快到正规医院妇科、泌尿外科或皮肤性病科就诊，
不要自行买药或拖延：

\begin{itemize}
  \item 疑似无保护行为后需要紧急避孕或暴露后预防（如 HIV 暴露）。
  \item 出现疼痛、异常出血、分泌物异常、溃疡、皮疹、发热等症状。
  \item 检测呈阳性，需要规范治疗（多数性传播感染可治愈或可有效控制）。
  \item 对避孕方法有疑问，或需要 HPV、乙肝等疫苗接种建议。
\end{itemize}

\chapter{性传播感染详解}
\label{chap:sti}

性传播感染（STI）是初学者最容易恐慌、也最容易被谣言误导的话题。
这一章用循证医学的视角把它讲清楚：哪些能治、哪些不能、如何发现、出了问题怎么办。

\section{三个基本事实}

\begin{enumerate}[leftmargin=2em,itemsep=0.3em]
  \item \textbf{多数 STI 没有症状。} 世界卫生组织指出，大多数性传播感染没有症状，
        女性以及咽部、直肠等非尿道部位尤其如此。所以"没感觉"不等于"没感染"，
        唯一确知方式是检测。
  \item \textbf{分为"可治愈"和"可控"两大类。} 细菌性和寄生虫类通常能根治；
        病毒类多为终身感染，但可用药物抑制，并不影响正常寿命。
  \item \textbf{早发现、早治疗能防并发症。} 未经治疗可导致不孕、宫外孕、死胎、
        神经系统与心血管疾病，并增加感染 HIV 的风险。
\end{enumerate}

\section{可治愈类：细菌与寄生虫}

\begin{table}[htbp]
  \centering
  \caption{常见的可治愈性传播感染}
  \begin{tabularx}{\textwidth}{@{}l X X@{}}
    \toprule
    \textbf{感染} & \textbf{常见表现} & \textbf{要点} \\
    \midrule
    衣原体 & 常无症状；可有分泌物增多、排尿灼痛 & 可治愈；不治可致女性生殖系统永久损伤、不孕或宫外孕 \\
    淋病 & 男性可有尿痛、黄绿色分泌物；女性常无症状或很轻微 & 可治愈，但耐药菌株增多，须完成全程，不可自行买药或分给别人 \\
    梅毒 & 一期为单个无痛溃疡；二期可全身皮疹（含手掌足底） & 可治愈；潜伏期无症状但血清学阳性 \\
    滴虫 & 女性分泌物异常、外阴瘙痒、性交痛 & 可治愈；性伴需同时治疗 \\
    \bottomrule
  \end{tabularx}
\end{table}

\noindent 共性：一旦确诊，\textbf{你和伴侣必须同时接受治疗}，否则容易反复交叉感染。

\section{不可治愈但可控类：病毒}

\begin{table}[htbp]
  \centering
  \caption{常见的病毒类性传播感染}
  \begin{tabularx}{\textwidth}{@{}l X X@{}}
    \toprule
    \textbf{感染} & \textbf{特点} & \textbf{应对} \\
    \midrule
    HIV & 攻击免疫细胞；感染头几个月可能无症状但传染性强 & 无法治愈，但抗逆转录病毒治疗（ART）可长期抑制，病毒量测不出时不会传染性伴 \\
    生殖器疱疹（HSV） & 疼痛性水疱或溃疡，易复发；很多人无症状或极轻微 & 终身感染，抗病毒药可减少复发；不会经马桶、泳池、毛巾等日常物品传播 \\
    乙肝（HBV） & 可经性接触传播 & \textbf{有疫苗}；慢性感染需长期管理 \\
    HPV & 极常见，约九成感染两年内自行消退；少数持续感染可致宫颈等癌症 & \textbf{有疫苗}；宫颈癌筛查（HPV 检测、TCT）可早发现癌前病变 \\
    \bottomrule
  \end{tabularx}
\end{table}

\noindent 关键区别：病毒类"治不好"不等于"没救"。ART 抑制 HIV、抗病毒药控制疱疹、
疫苗预防乙肝和 HPV，都是成熟手段。

\section{症状与窗口期}

出现以下任一信号应尽快就医：异常分泌物、生殖器溃疡或肿物、排尿灼痛、下腹痛、
性交痛、发热皮疹。但更要记住——\textbf{没有任何信号也可能已经感染}。

以 HIV 为例，多数人在感染后 28 天内产生抗体；WHO 建议高危暴露后检测阴性者，
\textbf{28 天后再复查一次}，因为过早检测可能出现假阴性。其他感染的检测时机各有不同，
请遵医嘱安排。

\section{检测：怎么做、多久一次}

\begin{itemize}[leftmargin=2em]
  \item \textbf{方式}：多为尿液或拭子采样（尿道、宫颈、咽部、直肠），无创快捷。
        主动告诉医生你的性行为方式，便于选对采样部位。
  \item \textbf{频率}：CDC 建议有性行为的男男性行为者至少每年查 HIV、梅毒、乙肝、丙肝（依风险）；
        25 岁以下或有新性伴、多性伴、性伴已感染的女性每年查衣原体与淋病。
        普通人确定固定伴侣前双方各查一次，是负责任的做法。
  \item \textbf{阳性不是终点}：多数 STI 可治愈或可控，规范治疗、通知伴侣、按期复查即可。
\end{itemize}

\section{暴露后怎么办：PEP}

HIV 暴露后预防（PEP）用于可能接触 HIV 的紧急情况（如安全套破裂、共用针具、性侵）：
\textbf{越早越好，最好 24 小时内、最晚不超过 72 小时}到急诊或医疗机构启动，
\textbf{每天服药、连服 28 天}。PEP 不能替代常规防护；若反复暴露，
应咨询医生改用暴露前预防（PrEP）。

在国内，各地疾控中心、传染病医院、三甲医院感染科或急诊可提供咨询与 PEP，
也可拨打当地 12320 卫生热线询问。

\section{预防小结}

\begin{enumerate}[leftmargin=2em,itemsep=0.3em]
  \item \textbf{安全套}：每次、全程、正确使用。对体液传播的感染（淋病、衣原体、HIV 等）效果好；
        对皮肤接触传播的感染（疱疹、HPV、梅毒）只能部分降低风险。
  \item \textbf{疫苗}：HPV、乙肝（及甲肝）疫苗。
  \item \textbf{固定性伴与双方检测}：双方均未感染的单一性伴关系风险最低。
  \item \textbf{定期筛查}：无症状感染只能靠筛查发现。
  \item \textbf{避免酒精与药物影响下决策}：会降低判断力，导致不戴套等风险行为。
\end{enumerate}

\chapter{关系与现实：小白最该知道的几件事}
\label{chap:reality}

前面各章偏重"知识"，这一章偏重"观念"。对初学者来说，
破除迷思、学会沟通、保护自己，比记住任何技巧都重要。

\section{破除对"第一次"的迷思}

\begin{itemize}[leftmargin=2em]
  \item \textbf{"出血 = 第一次"是过时观念}：阴道瓣（旧称处女膜）形态各异，
        很多人第一次并不出血；剧烈运动、卫生棉条等都可能影响它。
        用它评判别人既不科学也不尊重。
  \item \textbf{疼痛不是必须的}：放松、沟通、充分的前戏与润滑能大幅减少不适。
        如果很痛就停下来——不是"不够好"，而是方法或时机需要调整。
  \item \textbf{别追求"完美表现"}：第一次更重要的是彼此信任和沟通，不是表演。
  \item \textbf{影视夸大了戏剧性}：现实往往平淡得多，甚至有点失望，这完全正常。
\end{itemize}

\section{别把色情内容当教材}

如果你对性的"首次印象"来自网络视频，请意识到\textbf{那是表演，不是教育}：

\begin{itemize}[leftmargin=2em]
  \item 时长、体型、反应、场景都经过设计和剪辑，和真实生活差距很大。
  \item 它几乎不涉及\textbf{安全、同意、情感}——而这恰恰是现实中最重要的三件事。
  \item 长期拿它当标准，容易对自己的身体和伴侣产生不必要的不满与焦虑。
\end{itemize}

\section{性欲与节奏：没有标准答案}

\begin{itemize}[leftmargin=2em]
  \item 每个人的性欲高低、唤起方式、节奏都不同，还会随年龄、压力、健康、情绪波动。
  \item 有人浪漫先行，有人身体先行——都正常。
  \item \textbf{差异不等于问题}。真正值得关注的信号是：欲望突然、显著地变化
        并带来困扰，那可以咨询医生或心理咨询师。
\end{itemize}

\section{学会说"不"，也学会接受"不"}

这是初学者最需要、却最少被教的一课：

\begin{itemize}[leftmargin=2em]
  \item \textbf{拒绝不需要理由}："我不想"就是完整理由，不必解释到对方满意。
  \item \textbf{具体说法}："我现在不想""我还没准备好""换个方式吧""我累了想休息"。
  \item \textbf{对方说"不"时}：尊重、不追问、不施压、不冷战。那不是拒绝你这个人。
  \item \textbf{危险信号}：用"如果你爱我你就会""大家都这样""你不配合我就生气"
        来施压——这是操控，不是亲密。
\end{itemize}

\section{身体意象与自我接纳}

\begin{itemize}[leftmargin=2em]
  \item 身体的形状、大小、肤色、体毛千差万别，\textbf{没有"标准款"}。
  \item 图片和视频里的光影、角度、剪辑都会改变呈现，真实的人不是那样。
  \item 好的亲密关系建立在\textbf{接纳}上，而不是把自己或伴侣和某个图像比较。
\end{itemize}

\section{几个容易被忽视的实际问题}

\begin{itemize}[leftmargin=2em]
  \item \textbf{酒后 / 药后}：判断力和自控力下降。提前想好界限，不单独赴不明场合的约。
  \item \textbf{线上交友}：不轻易透露住址、单位、经济状况；初次见面选公共场所并告知朋友；
        \textbf{不发送私密照片}——一旦发出难以收回，可能被用于敲诈。
  \item \textbf{隐私}：任何亲密行为都属私密。\textbf{未经同意拍摄或传播他人私密影像，
        是严重伤害，且违法。}
  \item \textbf{金钱与礼物往来}：保持清醒，别被物质裹挟做出不愿的决定。
\end{itemize}

\section{如果感到困惑或受伤}

\begin{table}[htbp]
  \centering
  \caption{遇到困惑或伤害时的求助渠道}
  \begin{tabularx}{\textwidth}{@{}l X@{}}
    \toprule
    \textbf{情况} & \textbf{求助渠道} \\
    \midrule
    性健康知识、检测咨询 & 正规医院、当地疾控、12320 卫生热线 \\
    情绪困扰、心理压力 & 心理咨询师、心理援助热线 \\
    性骚扰、性侵 & 立即离开现场、拨打 110、保留证据、尽快就医并联系可信任的人 \\
    \bottomrule
  \end{tabularx}
\end{table}

\noindent\textbf{一句话总结}：慢一点、问清楚、尊重彼此、保护好自己。
任何时刻觉得不确定或不舒服，都可以停下来——查证或求助，一点都不丢人。

\chapter{见面了怎么办：从邀约到在场的清单}
\label{chap:meeting}

"见面"这件事，焦虑往往来自不确定。把流程拆开，每一步都知道该做什么，就不慌了。
本章按时间顺序分为三个阶段，并给出各自的红线。

\section{见面前：先把安全底线定好}

\begin{itemize}[leftmargin=2em]
  \item \textbf{选对场合}：第一次见面选\textbf{公共场所}（咖啡馆、商场、公园），
        时间尽量在白天或傍晚，不去对方指定的私密场所或陌生地点。
  \item \textbf{自己来、自己走}：自己安排交通，不搭对方的车去偏僻处；
        告知朋友或家人去哪、和谁、大概几点回。
  \item \textbf{管好饮品}：饮料或杯子不离视线；离开过又回来的不再喝；不接受陌生人递来的酒。
  \item \textbf{提前想清楚边界}：去之前就问自己——能接受到哪一步？什么情况一定喊停？
        把答案在心里定好，现场就不容易被带着走。
  \item \textbf{信任直觉}：如果对方催你换地方、要求保密、或不尊重你说"不"，这些都是信号。
\end{itemize}

\section{见面中：感受与沟通}

\begin{itemize}[leftmargin=2em]
  \item \textbf{先感受自己}：放松吗？还是隐隐不安？身体紧绷、想离开的念头，
        都是重要信号，别忽略。
  \item \textbf{边界说在前面}：在气氛还好时自然地说"我今天想先聊聊天""我不去太晚的地方"。
        提前说比临场拒绝容易得多。
  \item \textbf{随时可以停}：任何时刻不舒服，直接说"我想回去了"即可，不必解释到对方满意。
  \item \textbf{观察对方对你说"不"的反应}：提一个小要求或小拒绝，看他是尊重还是施压、冷脸、纠缠——
        这能看出这个人值不值得继续。
  \item \textbf{别被酒精带节奏}：微醺会降低判断力，保持清醒才能真正地同意或拒绝。
\end{itemize}

\section{想继续亲密：现场的四个确认}

\begin{enumerate}[leftmargin=2em,itemsep=0.3em]
  \item \textbf{确认同意}：双方都清醒、自愿、能随时喊停。"不抗拒"不等于"同意"，
        需要明确的肯定。
  \item \textbf{安全套准备好}：自己备好、在有效期内、知道怎么用，别指望对方一定有。
  \item \textbf{谈避孕与检测}：可以自然地说"我们用安全套好吗""你最近做过检测吗"——
        这不是扫兴，是负责。
  \item \textbf{慢一点、多沟通}：节奏由双方舒服的程度决定；疼、不适、犹豫就停。
        第一次不理想很正常，不必有压力。
\end{enumerate}

\noindent\textbf{三条红线（碰到就走）}：无保护也无所谓、对方不情愿还硬来、
用感情或压力逼你改变决定。任何一条出现，立刻停止、离开、联系朋友。

\section{见面之后}

\begin{itemize}[leftmargin=2em]
  \item \textbf{沟通感受}：事后可简单聊聊彼此的感受，好的和不好的都可以说。
  \item \textbf{留意身体}：如有无保护行为，记住时间窗——紧急避孕越早越好，
        HIV 暴露后预防（PEP）72 小时内；需要时按第~\ref{chap:sti}~章安排检测。
  \item \textbf{保护隐私}：若对方拍摄私密影像，未经你同意不得传播；
        遇到胁迫，第一时间求助（110 / 12320 / 可信任的人）。
  \item \textbf{允许自己复盘}：不管结果如何，想想"哪些我做对了、哪些下次可以不同"，这就是成长。
\end{itemize}

\noindent\textbf{一句话总结}：见面的核心不是"表现好"，而是\textbf{安全 + 自愿 + 能随时喊停}。
守住这三点，其余的慢慢来就好。

\chapter{男性视角的补充：同样重要的一半}
\label{chap:male}

前面各章更多是"如何保护自己"的口吻。但性是两个人的事，
男性有对等的权利、责任与困扰。这一章从男性角度补齐。

\section{同意是双向的：你也有权说不}

\begin{itemize}[leftmargin=2em]
  \item 前面讲"学会说不"，对男性同样成立。\textbf{对方有同意权，你也有}。
  \item 不必因为"是男人"就必须主动、必须想、必须持续。没准备好就说"我今天不想"，不必解释。
  \item 被对方施压（"你不配合就是不爱我"）和你施压一样，都是越界。
  \item \textbf{别把"不抗拒"当成"同意"}——明确、清醒、可撤回的肯定，才是同意。
\end{itemize}

\section{男性也会成为受害者}

这不是小事。WHO 数据显示，约\textbf{1/5 的女性、1/7 的男性}报告童年遭受过性侵。
男性遭遇性骚扰、性侵、亲密关系暴力的比例不低，但因为"男人就该扛着"的观念，
很多人选择沉默。

\begin{itemize}[leftmargin=2em]
  \item 遭遇侵害不是你的错，与"强不强壮"无关。
  \item 求助渠道对男性同样开放：拨打 110、12320，或联系医院、可信任的人。
  \item WHO 明确指出：僵化的"男性气概"规范会损害男性健康——鼓励冒险、酗酒、不愿求助。
        主动求助是勇敢，不是软弱。
\end{itemize}

\section{常见身体困扰：常见、可治疗、别乱补}

\subsection*{1. 勃起功能障碍（ED）是全身健康的信号}

\begin{itemize}[leftmargin=2em]
  \item 不同研究数字差异很大（如波士顿 MMAS 研究 40--70 岁约 52\%，
        科隆研究 30--80 岁约 19.2\%），共同结论是\textbf{随年龄上升}。
  \item 关键：ED 常与糖尿病、高血压、血脂异常等代谢/心血管问题相关。
        EAU 指南要求在诊治时评估血糖、血脂和睾酮——\textbf{它可能是身体给出的早期警报}，
        值得正规就医，而不是买补药。
  \item 治疗是一套成熟方案：先改善生活方式（运动、减重、控压控糖），
        一线药物为 PDE5 抑制剂，均需在医生指导下使用。
\end{itemize}

\subsection*{2. 早泄（PE）：先搞清楚是不是"病"}

\begin{itemize}[leftmargin=2em]
  \item ICD-11 已把它改称"早期射精"：指插入前或插入后极短时间内射精、
        几乎无控制感、持续数月并带来明显困扰。
  \item EAU 强调\textbf{不能只看时间}，还要看控制感、是否普遍、是否造成困扰。
        "正常变异"（时快时慢）和"主观性早泄"（时间其实正常但自认为快）都不算真正的疾病。
  \item 真正需要时，行为训练、盆底肌锻炼、药物等都是有效手段，去泌尿外科/男科评估即可。
\end{itemize}

\subsection*{3. 别信"壮阳保健品"}

来路不明的补药、偏方成分不明，可能有害且耽误真正病因的诊治。有困扰就去看正规医院。

\section{男性在避孕与防病中的责任}

\begin{itemize}[leftmargin=2em]
  \item \textbf{安全套}：男性主导使用的工具，主动准备好、正确佩戴，是责任也是保护自己。
  \item \textbf{输精管结扎}：高效、永久性的男性避孕方式。CDC 数据：典型使用第一年失败率约
        0.15\%；术后 8--16 周需做精液检查确认无精子，确认前仍需用屏障避孕。
        它\textbf{不预防性病}。后悔率约 5\%，决定前要充分咨询。
  \item \textbf{HPV 疫苗男性也要打}：CDC 建议男女都在 11--12 岁常规接种、可 9 岁起，
        未接种者 26 岁前补种，27--45 岁与医生共同决策。可预防生殖器疣、肛门癌等。
  \item \textbf{乙肝疫苗}：有多个性伴或接受性病诊疗者应接种。
  \item 定期做性病筛查，同样是对自己和伴侣负责。
\end{itemize}

\section{"表现焦虑"多半是被教出来的}

\begin{itemize}[leftmargin=2em]
  \item 色情内容里的时长、强度、反应都是表演，拿它当标准必然焦虑。
  \item 一次性生活"不理想"非常常见——疲劳、压力、饮酒、情绪都会影响。
        偶尔失常不等于"不行了"。
  \item 越关注"表现"，越容易真的出问题（焦虑本身就是 ED 的常见原因）。
        把注意力放回连接、沟通和感受，反而更好。
  \item 真正需要警惕的是问题\textbf{持续出现}并带来困扰——那就去就医，而不是独自硬扛。
\end{itemize}

\section{给男性的一页清单}

\begin{enumerate}[leftmargin=2em,itemsep=0.3em]
  \item 我尊重对方的同意，也保有自己说"不"的权利。
  \item 我知道男性也可能成为受害者，遇到侵害会求助而不是硬扛。
  \item 我明白 ED/早泄常见且可治疗，ED 还可能是全身健康信号，会去正规医院而非买补药。
  \item 我会主动准备安全套，知道输精管结扎与 HPV 疫苗等选项。
  \item 我不拿色情内容当标准，不把一次性失常当成"不行"。
  \item 我不迷信"男人气概"，会表达情绪、会求助。
\end{enumerate}

\chapter{到了私密场所后：以酒店为例的具体清单}
\label{chap:private}

进了房间\textbf{不等于同意任何事}。环境变了，权利一点没变——同意依然必须清醒、
自愿、可随时撤回。而且这种场景封闭、隐蔽、退路少，反而更要主动把安全底线再确认一遍。
下面按时间顺序拆成四步：进房间 $\to$ 定边界 $\to$ 过程中 $\to$ 之后。

\section{第一步：进房间，先花两分钟看环境}

\begin{itemize}[leftmargin=2em]
  \item 选正规、口碑好的酒店，\textbf{自己登记、自己开门}；进门先试门锁能否正常反锁、
        防盗扣是否完好。
  \item 给朋友或家人发个房号/定位；\textbf{手机充满电}，放在伸手可及的地方。
  \item 别急着放松，先做一次快速隐私排查（见下）。
\end{itemize}

\subsection*{关于偷拍：先纠正一个流传很广的误区}

网上说"关灯用手机摄像头扫，有红点就是针孔摄像头"——\textbf{这个方法基本不靠谱}。
科技日报、中国新闻网均指出：绝大多数针孔摄像头因体积限制没有夜视红外功能，不会发红光；
而现在的手机镜头大多带滤光片，即便有红外光也显示不出来。

更实用的笨办法：

\begin{enumerate}[leftmargin=2em,itemsep=0.3em]
  \item \textbf{手电筒贴近看小孔}：关灯，用手电近距离照——插座面板、台灯底座、
        空调出风口、烟雾报警器、电视底座、镜子边缘与背面、床头布艺缝隙、门锁钥匙孔。
        镜头是玻璃制品，光照下会有\textbf{反光点}。
  \item \textbf{查 WiFi 列表}：连上房间 WiFi，留意 \texttt{IPC}、\texttt{Camera}、
        \texttt{DV} 之类陌生设备名。
  \item \textbf{物理遮挡最稳妥}：用浴巾/衣物盖住正对床的可疑物件，拉窗帘、关灯，
        让它"看不见、拍不清"。
  \item \textbf{发现疑似装置}：先拍照录像留存证据，再找酒店交涉，必要时报警。
        法律上，宾馆房间属于私密空间（《民法典》第 1033 条），偷窥、偷拍、散布他人隐私
        可按《治安管理处罚法》罚款拘留；非法生产、销售窃照专用器材则触及《刑法》第 283 条。
\end{enumerate}

\begin{quote}
这类事件毕竟是个案，简单自查是为多一份安心，不必过度紧张。
\end{quote}

\section{第二步：把"同意与边界"当面再说一遍}

很多误会源于"我以为你同意了"。主动说清楚，反而最稳：

\begin{itemize}[leftmargin=2em]
  \item \textbf{再确认意愿}："我们确定要这样吗？"——你自己和对方的 YES 都必须明确、清醒。
        \textbf{沉默、不抗拒、之前答应过，都不等于此刻的同意。}
  \item \textbf{再确认防护}："我们用安全套好吗？"——自己备好、在有效期内、知道怎么用，
        别指望对方一定有。
  \item \textbf{说清边界}：哪里可以、哪里不行、什么情况要停。提前说，比临场拒绝容易得多。
  \item \textbf{约定一个"停"字}：说了就立刻停，不问原因。
\end{itemize}

\section{第三步：过程中，守住四条线}

\begin{enumerate}[leftmargin=2em,itemsep=0.3em]
  \item \textbf{全程清醒}：饮料别离开视线，离开过又回来的不再喝。判断力下降时，
        谈不上真正的同意。
  \item \textbf{安全套全程、正确用}：接触前就戴；只用\textbf{水基或硅基}润滑剂，
        油基会削弱乳胶。
  \item \textbf{随时可以停}：不舒服、犹豫、疼，直接说停——不需要理由，也不需要对方满意。
  \item \textbf{警惕"偷偷摘套"（stealthing）}：未经同意在过程中摘掉安全套，
        是对同意的严重违背。研究显示遭遇者多数有明显情绪压力，澳大利亚、加拿大、德国、
        丹麦等多国已将其立法追责。一旦发现：\textbf{立刻停止，要求对方重新戴上；
        不照做就离开}。这不是小题大做，是你的健康与身体自主权。
\end{enumerate}

\begin{quote}
还有一条：\textbf{未经你同意拍摄或传播亲密影像，是严重伤害，且违法。}
无论对方说"就拍着玩"还是"马上删"，你都有权拒绝、并要求当场删除。
\end{quote}

\section{第四步：之后，别忽略收尾}

\begin{itemize}[leftmargin=2em]
  \item \textbf{若无保护、安全套破裂或滑脱}：记住两个时间窗——\textbf{紧急避孕通常 72 小时内}
        （部分方案可到 120 小时，越早越好，不作常规手段）；
        \textbf{HIV 暴露后预防（PEP）最晚 72 小时内}启动，连服 28 天。
        尽快就医或咨询药师，并按第~\ref{chap:sti}~章安排后续检测。
  \item \textbf{隐私}：拍过的影像确认删除；不留存、不转发任何私密内容。
  \item \textbf{沟通与复盘}：聊聊彼此感受；想想哪些做对了、哪些下次可以不同。
  \item \textbf{若遭遇侵害}：立即离开、保留证据、拨打 110 或 12320、尽快就医并联系
        可信任的人。\textbf{这不是你的错。}
\end{itemize}

\section{一张随身清单}

\begin{enumerate}[leftmargin=2em,itemsep=0.3em]
  \item 正规酒店、自己登记、门锁反锁、行踪告知朋友。
  \item 两分钟查偷拍：手电看小孔反光 + WiFi 陌生设备 + 浴巾遮挡。
  \item 再确认：我愿意、对方愿意、都用安全套、说清边界。
  \item 全程清醒、饮品不离视线、安全套全程正确用。
  \item 任何不适立刻停；发现偷偷摘套/偷拍，立即止损离开。
  \item 记牢 72 小时：紧急避孕 / PEP。
  \item 私密影像不拍不传；遇侵害 110 / 12320 / 就医。
\end{enumerate}

\chapter{床上的事：亲密过程的健康框架}
\label{chap:intimacy}

本章不是"技巧手册"，也不涉及露骨的动作描写——那不是健康科普该做的事，
也与全书"安全、自愿、知情"的基调不符。但"整个过程按什么顺序推进、哪里最容易出问题、
怎样让双方都舒服"是完全正当、也最能让初学者心里有底的内容。
核心四个字：\textbf{慢、沟通、全程防护、随时可停}。

\section{开始前：把氛围和心态摆正}

\begin{itemize}[leftmargin=2em]
  \item \textbf{不表演、不赶时间}。第一次紧张、笨拙、笑场都很正常。把注意力放在彼此的连接和感受上，
        而非"表现好不好"。
  \item \textbf{环境}：干净、舒适、温度适宜；安全套和水基润滑剂放在伸手可及处，免得到时手忙脚乱。
  \item \textbf{同意是持续的，不是一次性的}。进了房间、开始亲密、每一个新动作前，
        都可以自然确认："这样可以吗？""你喜欢这样吗？"——对方明确、清醒的回应才是继续的前提。
\end{itemize}

\section{前戏：不是"例行公事"，是让身体做好准备}

前戏（拥抱、亲吻、爱抚、语言交流）的目的是让双方放松、唤起、建立信任，而非机械走流程：

\begin{itemize}[leftmargin=2em]
  \item \textbf{唤起需要时间，个体差异很大}。尤其女性往往需要更长时间；不要急着进入"正题"。
  \item \textbf{多问、多观察}："这样舒服吗？""要不要慢一点？"对方的呼吸、表情、身体松紧都是重要信号。
  \item \textbf{润滑不足是疼痛和不适的常见原因}。感到干涩就补\textbf{水基或硅基}润滑剂
        （油基会削弱乳胶安全套）——这不是"不够爱"的表现，而是正常生理差异，
        年龄、压力、药物都会影响。
\end{itemize}

\section{安全套：全程、正确、别侥幸}

CDC 强调安全套必须\textbf{每次性行为都全程使用}才有效。结合场景复习最易错的几点：

\begin{enumerate}[leftmargin=2em,itemsep=0.3em]
  \item 在生殖器接触\textbf{之前}就戴好：勃起后捏住储精囊排出空气、展开到根部。
  \item 润滑只用水基/硅基；凡士林、婴儿油、护肤乳等\textbf{油基会让乳胶破裂}。
  \item 射精后按住套口、趁勃起退出；用后打结丢弃，不可重复使用。
  \item \textbf{中途破裂或滑脱}：立即停止、换新；按第~\ref{chap:sti}~章与紧急避孕的时间窗处理（72 小时内）。
  \item 对方说"不用套""就这一次""我去检测了没事"——\textbf{坚持是你的权利}，任何一方都有一票否决权。
\end{enumerate}

\section{过程中：慢、沟通、看反应}

没有"标准剧本"，但有几条通用安全原则：

\begin{itemize}[leftmargin=2em]
  \item \textbf{循序渐进}，不急躁、不勉强；任何新动作先看对方反应、听对方意愿。
  \item \textbf{持续沟通}：一句"还好吗？""要慢一点吗？"比沉默猜测可靠得多。
  \item \textbf{疼痛、干涩、不适、犹豫}——立刻停下来，问清楚、调整或改天。
        \textbf{疼不是"必经之路"，是身体在提醒你。}
  \item 保持清醒。酒精会降低判断力，也会让"同意"变得可疑。
\end{itemize}

\section{关于"高潮"：放下 KPI}

\begin{itemize}[leftmargin=2em]
  \item \textbf{高潮不是必须的}，也不是衡量成功与否的指标。WHO 强调良好的性健康是
        "积极、尊重的体验，可以带来愉悦且安全"，重点在连接与感受，不在某个生理终点。
  \item 女性是否高潮、多久高潮，个体差异极大；很多人需要特定刺激，也有人并不在意——
        \textbf{都正常}，不必焦虑或假装。
  \item 一次性生活不理想很常见：疲劳、压力、饮酒、情绪都会影响。偶尔失常不等于"不行了"。
  \item 真正需要警惕的是问题\textbf{持续出现并带来困扰}——那时去正规医院（泌尿外科/男科/妇科）
        或咨询专业人士，而不是自己硬扛或乱补。
\end{itemize}

\section{事后：别忽略这段"隐形"的重要时光}

\begin{enumerate}[leftmargin=2em,itemsep=0.3em]
  \item \textbf{温存与沟通}：拥抱、聊聊彼此的感受——好的、不好的都可以说。这能加深信任，比任何技巧都管用。
  \item \textbf{排尿}：事后尽快排尿，有助于冲掉进入尿道的细菌、降低尿路感染风险。
        CDC 明确：性行为会把细菌带入尿道，女性因尿道更短而更易感染。
        若出现尿频、尿急、灼痛、下腹痛，及时就医。
  \item \textbf{清洁}：用清水适度清洗外阴即可，不必过度使用洗液；不共用毛巾等贴身物品。
  \item \textbf{若发生无保护行为}：记住 72 小时——紧急避孕（越早越好，不作常规手段）、
        HIV 暴露后预防（PEP，最晚 72 小时内启动、连服 28 天），并按第~\ref{chap:sti}~章安排检测。
  \item \textbf{隐私}：不拍摄、不留存、不传播任何亲密影像；对方提出拍摄，你有权拒绝并要求删除。
\end{enumerate}

\section{任何时候都可以停：四条红线}

出现以下任一情况，\textbf{立刻停止、离开、求助}：

\begin{enumerate}[leftmargin=2em,itemsep=0.3em]
  \item 任何一方说"停"、犹豫、沉默或明显不舒服；
  \item 疼痛或身体不适；
  \item 对方偷偷摘掉安全套（stealthing）或坚持无保护；
  \item 未经同意的拍摄、胁迫，或任何让你觉得不安全的行为。
\end{enumerate}

求助渠道：\textbf{110}（人身安全）、\textbf{12320}（卫生热线）、可信任的家人朋友、正规医院。
这不是你的错，也不需要"再忍一忍"。

\chapter{亲密过程中的身体反应：性反应周期与注意事项}

\begin{center}
\begin{tikzpicture}[font=\small]
  \draw[mygray,thin] (0,0) -- (9.5,0) node[midway,below,mygray]{时间 →};
  \draw[mygray,thin] (0,0) -- (0,3.6) node[midway,left,mygray]{唤起水平};
  \draw[myblue,thick] (0,0.4) to[out=45,in=170] (2.2,3.0) to[out=0,in=180] (4.6,2.7) to[out=0,in=110] (5.6,3.4) to[out=350,in=90] (9,0.5);
  \fill[myblue] (5.6,3.4) circle (0.06);
  \node[mygray] at (1.1,3.9){兴奋期};
  \node[mygray] at (3.4,3.9){平台期};
  \node[myblue] at (5.6,3.9){高潮};
  \node[mygray] at (7.4,3.9){消退期};
  \node[mygray,font=\footnotesize] at (1.1,-0.5){唤起·充血·润滑};
  \node[mygray,font=\footnotesize] at (3.4,-0.5){维持高位·张力升};
  \node[myblue,font=\footnotesize] at (5.6,-0.5){节律性收缩};
  \node[mygray,font=\footnotesize] at (7.6,-0.5){回落基线};
\end{tikzpicture}
\par\vspace{0.5em}\noindent\footnotesize 图：性反应周期（Masters \& Johnson）四阶段唤起水平曲线。
\end{center}
\label{chap:cycle}

本章讲"过程中身体到底会发生什么"——用经典的\textbf{性反应周期}四阶段模型，
从生理角度把过程说清楚。全程用医学科普的口吻，聚焦正常反应、个体差异和注意事项，
不涉及露骨的动作描写。先记住一句：这些反应大多是自动发生的，你不需要"表演"，
也不需要拿自己去对照任何标准。

\section{整体框架：四个阶段}

经典模型（Masters \& Johnson，1966 年《Human Sexual Response》）把一次亲密过程的身体反应
分为四个阶段：\textbf{兴奋期 $\to$ 平台期 $\to$ 高潮期 $\to$ 消退期}。
它不是机械流水线——各阶段长短因人而异，有时会跳过或反复，情绪、疲劳、压力、环境都会影响。
理解它，主要是为了知道"什么反应是正常的""什么时候该注意"。

\section{兴奋期：身体开始"准备"}

\textbf{身体会发生什么}（男女共同）：心跳加快、呼吸变快、血压上升、肌肉逐渐紧张、皮肤可能泛红；
男性阴茎勃起；女性阴道分泌润滑液、阴道扩张、外生殖器充血。

\begin{itemize}[leftmargin=2em]
  \item \textbf{唤起需要时间}，个体差异很大。女性往往需要更长时间，不要急着进入"正题"——这也是前戏的意义。
  \item \textbf{勃起受情绪和状态影响很正常}。紧张、疲劳、饮酒、压力大时勃起不理想很常见，不等于"有问题"。
  \item \textbf{干涩就补水基润滑剂}，这是正常生理差异（年龄、压力、药物都会影响），不是"不够投入"。
  \item 保持沟通："这样舒服吗？"比默默猜测可靠。
\end{itemize}

\section{平台期：反应持续加强}

\textbf{身体会发生什么}：心率、呼吸进一步加快，肌肉紧张加剧，部分人胸腹部出现红疹（"性红晕"）；
男性尿道口可能分泌少量液体；女性阴道下段收缩、相关组织进一步充血。

\begin{itemize}[leftmargin=2em]
  \item 这是"持续"阶段，节奏可根据双方感受调整——\textbf{快了就慢一点，累了就停一停}。
  \item 男性注意：\textbf{射精前的分泌物中可能已混有精子}，所以"体外排精"并不可靠，
        不能当作避孕方法（见第~\ref{chap:intimacy}~章）。
  \item 若一直"卡"在此阶段、许久没有进展，不必焦虑——放慢、换方式或顺其自然结束都可以。
        \textbf{亲密不等于必须走到某一步。}
\end{itemize}

\section{高潮期：短暂的峰值}

\textbf{身体会发生什么}：盆底及相关肌肉出现节律性收缩，心率、呼吸达到高峰，
主观上有强烈愉悦感；持续几秒到几十秒。男性射精通常与高潮同时发生，
女性是否高潮、强弱差异极大。

\begin{itemize}[leftmargin=2em]
  \item \textbf{女性高潮差异非常大}：有人容易、有人需要特定条件、也有人从未体验过——
        \textbf{都正常}，不必假装，更不要用"对方有没有高潮"评判这次亲密是否成功。
  \item 高潮不是必须的终点。\textbf{连接、愉悦、彼此舒服才是目的}，高峰只是可能的组成部分。
  \item 一次不理想很正常，别给自己下"不行"的定义。
\end{itemize}

\section{消退期：身体慢慢恢复}

\textbf{身体会发生什么}：心跳、呼吸、肌肉紧张在数分钟内逐渐恢复。
男性通常有\textbf{不应期}（refractory period）——此期间难以再次勃起，年轻时较短、随年龄增长变长；
女性一般没有明显不应期，但也可能因敏感而需要休息。

\begin{itemize}[leftmargin=2em]
  \item \textbf{不应期是正常的生理现象}，不是"能力下降"。别给自己或对方"马上再来"的压力。
  \item 消退期是\textbf{温存的好时机}：拥抱、说说话、聊聊感受。这段"事后时光"对关系的意义，
        往往比高峰本身更大。
  \item 若有无保护行为，记住 72 小时窗口（紧急避孕 / PEP），并按第~\ref{chap:sti}~章安排检测。
\end{itemize}

\section{贯穿全程的注意事项}

\begin{enumerate}[leftmargin=2em,itemsep=0.3em]
  \item \textbf{没有"标准"}：时长、反应、节奏人人不同，也随年龄、状态、情绪变化。
        别拿自己或伴侣和任何图像、视频比较。
  \item \textbf{安全套全程、正确使用}：这是全程不变的底线（用法见第~\ref{chap:intimacy}~章）。
  \item \textbf{保持清醒}：酒精会降低判断力，也让"同意"变得可疑。
  \item \textbf{持续沟通}：任何新动作前确认意愿，"还好吗？""要不要慢一点？"——沉默不是同意。
  \item \textbf{疼痛、出血、明显不适就停}，并考虑就医，不要忍。
  \item \textbf{情绪连接高于一切}：WHO 把良好的性健康定义为"积极、尊重、愉悦且安全的体验，
        免于胁迫、歧视和暴力"。技术是次要的，尊重与舒服才是核心。
\end{enumerate}

\section{几个常见困惑}

\begin{itemize}[leftmargin=2em]
  \item \textbf{"没有高潮是不是有问题？"}——多数情况下不是。关注愉悦和连接，而非某个生理终点；
        若长期困扰可咨询医生。
  \item \textbf{"中途软了/勃起不理想"}——偶发很常见；若持续出现并带来困扰，
        去正规医院（泌尿外科/男科）评估，别乱吃补药。
  \item \textbf{"总觉得太快了"}——先分清是不是"正常变异"或"主观感受"，别急着贴标签；
        持续困扰可就医（EAU 指南强调不能只看时间）。
  \item \textbf{"总觉得疼/干"}——先试水基润滑、放慢节奏；若持续疼痛，就医排查原因。
\end{itemize}

\chapter{亲密行为的类型与健康防护}
\label{chap:practices}

本章属于全面性教育里"行为方式与健康"的正当部分：用客观语言讲清楚有哪些常见行为、
每种的健康关注点、以及该怎么防护。目的是科学认知和安全，不是"技巧指南"，也不是情色描写。
内容依据 WHO、CDC 的循证资料。

\section{前提：任何行为都要满足的四个条件}

\begin{itemize}[leftmargin=2em]
  \item \textbf{双方自愿、明确同意}——不是"不抗拒"，是清楚的 YES。
  \item \textbf{双方清醒}——酒精、药物会影响判断，也让"同意"失效。
  \item \textbf{有相应防护}——按行为类型选对屏障工具（见下）。
  \item \textbf{随时可停}——任何不适，一句话就停，不问原因。
\end{itemize}

任何一条不满足，这件事就不该发生。这一条比后面所有内容都重要。

\section{常见行为的类型与健康关注点}

\begin{center}
\begin{tabularx}{\linewidth}{@{}l X X@{}}
\toprule
\textbf{类型} & \textbf{主要健康关注点} & \textbf{推荐防护} \\
\midrule
非插入类（拥抱、亲吻、爱抚、相互自慰）
  & 皮肤接触传播的感染、手部卫生、共用器具
  & 勤洗手；不共用未清洁器具；避免体液接触黏膜或伤口 \\
\addlinespace
口交（口--生殖器/肛门）
  & 淋病、衣原体、梅毒、疱疹、HPV；口--肛时的肠道菌群
  & 对阴茎用安全套；对阴道/肛门用口交膜 \\
\addlinespace
阴道性交
  & 意外怀孕 + 体液传播的 STI
  & 安全套每次全程 + 水基/硅基润滑 \\
\addlinespace
肛交
  & 黏膜更易微小损伤；直肠没有自然润滑
  & 安全套 + 充足水基/硅基润滑；换套后再接触其他部位 \\
\bottomrule
\end{tabularx}
\end{center}

关键认知（CDC/WHO 反复强调）：

\begin{itemize}[leftmargin=2em]
  \item \textbf{感染可以发生在任何部位}——口腔、肛门、直肠都可能被感染，而且常常没有症状。
        "看起来健康"说明不了什么。
  \item \textbf{安全套能显著降低体液传播的感染}（淋病、衣原体、HIV 等），
        但对\textbf{皮肤接触传播的}（疱疹、梅毒、HPV）保护有限——未覆盖的皮肤仍可能接触。
  \item \textbf{安全套是目前唯一能同时避孕和防病的工具}（WHO）。
\end{itemize}

\section{屏障工具怎么用}

\subsection*{安全套}

\begin{itemize}[leftmargin=2em]
  \item 每一次、每一段行为都用新的；阴道、肛门、口腔行为\textbf{各用一个新的}。
  \item 接触前就戴好；用\textbf{水基或硅基}润滑。凡士林、婴儿油、护肤乳、食用油等
        \textbf{油基会破坏乳胶}。
  \item 中途破裂或滑脱立即停止、换新；射精后按住套口趁勃起退出。
  \item 乳胶过敏者用聚氨酯等合成材质。
\end{itemize}

\subsection*{口交膜（dental dam）}

\begin{itemize}[leftmargin=2em]
  \item 用于口对阴道或口对肛门：乳胶/聚氨酯薄膜，铺在口与身体部位之间。
  \item \textbf{每次都用新的}，检查有效期与破损；开始前放好、结束前不移除。
  \item 可用\textbf{水基/硅基}润滑防破裂；不要拉伸、不要重复使用。
  \item \textbf{不要用 nonoxynol-9（杀精剂）}——它会刺激黏膜、反而增加风险。
  \item 临时没有口交膜，可用\textbf{安全套剪开成方形}替代（CDC 做法）。
  \item 用后丢弃垃圾桶，不冲马桶。
\end{itemize}

\subsection*{换套与清洁（常被忽略，但很关键）}

\begin{itemize}[leftmargin=2em]
  \item \textbf{肛门相关行为后，接触其他部位前换新安全套}，避免把肠道细菌带到尿道、阴道或口腔。
  \item 手部在接触不同部位前后\textbf{用肥皂水清洗}；涉及手指--肛门接触时 CDC 建议可戴一次性手套。
  \item 避免可能导致口腔接触粪便的行为（如口--肛直接接触），以降低肠道感染风险。
\end{itemize}

\section{为什么"多重防护"才叫合理}

单靠一样工具不够，因为不同感染传播途径不同：

\begin{itemize}[leftmargin=2em]
  \item \textbf{安全套}对体液传播（HIV、淋病、衣原体、乙肝）效果好；
  \item 但对\textbf{皮肤接触传播}（疱疹、梅毒、HPV）保护有限——未覆盖的皮肤仍可能接触；
  \item 因此 WHO/CDC 建议组合：\textbf{安全套 + 疫苗（HPV、乙肝）+ 双方检测 + 固定性伴 + 定期筛查}。
  \item \textbf{多数 STI 无症状}，"看起来健康"说明不了任何问题，唯一确知方式是检测
        （见第~\ref{chap:sti}~章）。
\end{itemize}

\section{关于"合理"的一条现实提醒}

如果脑中的"动作内容"来自成人影片，请把它清零重来：

\begin{itemize}[leftmargin=2em]
  \item 那是\textbf{表演}，时长、强度、反应都经过设计和剪辑，与真实亲密生活差距很大；
  \item 它几乎不涉及\textbf{同意、防护、沟通}——而这恰恰是现实中最重要的三件事；
  \item 拿它当模板，容易对伴侣和自己产生不切实际的期待与焦虑。
\end{itemize}

\textbf{合理的标准是}：双方清醒自愿、有防护、彼此舒服、随时能停、事后无担忧。
满足这些，就是好的；不满足，就停下来补上。

\section{什么时候检测/就医}

\begin{itemize}[leftmargin=2em]
  \item 无保护行为、安全套破裂/滑脱、或不确定是否暴露 $\to$ 记住 \textbf{72 小时}
        （紧急避孕 / HIV PEP），尽快就医；
  \item 按时间窗做 STI 筛查（过早可能假阴性，遵医嘱复查）；
  \item 出现分泌物异常、溃疡、排尿灼痛、下腹痛、发热皮疹 $\to$ 及时就医；
  \item 对屏障用法、疫苗、避孕有疑问 $\to$ 咨询医生、药师或 12320。
\end{itemize}

\chapter{非插入类行为详解：完整、合理、不遗漏}
\label{chap:outer}

这一章把"非插入类"单独拿出来讲清楚。很多人对它的认知停留在"前戏"，
但它的健康特征、防护方式和它作为完整亲密行为的价值，都值得单独了解。

\section{先纠正一个观念}

"只有插入才算性"是个误区。\textbf{非插入类}（英文常称 outercourse / edge sex）
本身就有完整的亲密价值和快感，对很多人来说，这甚至就是他们偏好或唯一想要的方式。
它也是关系里建立信任、探索彼此感受的低压力入口——不该被当成"前戏的附属"，
它就是完整的亲密行为。

\section{它具体包括哪些}

\begin{itemize}[leftmargin=2em]
  \item 拥抱、依偎、身体贴合；
  \item 亲吻（浅吻、深吻/法式吻）；
  \item 抚摸、爱抚（全身触碰、抚摸胸部与敏感部位）；
  \item 手刺激（用手刺激对方生殖器，男女均适用）；
  \item 摩擦（生殖器不插入的贴合摩擦）；
  \item 边缘性行为（接近高潮但不进入的状态）。
\end{itemize}

\section{健康要点：不是"零风险"，但风险特征不同}

这是最容易被误解的一点。关键事实（CDC）：

\begin{itemize}[leftmargin=2em]
  \item \textbf{HIV}：CDC 明确"接吻、相互手刺激"属\textbf{可忽略风险}，
        nPEP（暴露后预防）不常规推荐——除非"可见血液 + 黏膜/破损皮肤接触"。
        所以普通情况下 HIV 风险极低。
  \item \textbf{但}：疱疹（HSV）、HPV、梅毒是\textbf{皮肤--皮肤接触}传播的。
        即使没有插入、没有体液交换，\textbf{皮肤对皮肤、口--生殖器接触都可能传播}。
  \item \textbf{唾液交换}：深吻理论上可传播 HSV-1（口唇疱疹）；
        HIV 经唾液风险极低，除非有可见血/牙龈出血 + 口腔破损。
\end{itemize}

\textbf{一句话记住：HIV 看体液，疱疹/HPV/梅毒看皮肤接触。}
非插入类主要要防的是后者，而不是 HIV。

\section{每一类怎么注意}

\subsection*{拥抱 / 依偎 / 浅吻}

\begin{itemize}[leftmargin=2em]
  \item 一般风险很低。
  \item 深吻：留意对方有无\textbf{活动性口唇疱疹}（嘴唇起水疱）；若对方正在发作，避免深吻。
  \item 有口腔溃疡、牙龈出血、嘴角裂伤时，避免深吻和口腔接触。
\end{itemize}

\subsection*{抚摸 / 爱抚}

\begin{itemize}[leftmargin=2em]
  \item 主要注意\textbf{手部卫生}：接触前洗手、指甲修平干净（避免刮伤黏膜，破损皮肤会提高传播风险）。
  \item 避免直接触摸对方的\textbf{活动性溃疡或疣体}（疱疹水疱、HPV 疣）。
\end{itemize}

\subsection*{手刺激}

\begin{itemize}[leftmargin=2em]
  \item 洗手、修剪指甲，动作轻柔，避免指甲划破黏膜。
  \item \textbf{关键}：若手在肛门/生殖器/口之间移动，每次换部位前\textbf{洗手或戴一次性手套}
        （CDC 对涉及肛门接触的建议）；避免把带体液的手直接放到自己或对方嘴里。
  \item 有体液接触后，事后清洗手部。
\end{itemize}

\subsection*{摩擦 / 边缘}

\begin{itemize}[leftmargin=2em]
  \item 主要关注皮肤--皮肤接触传播疱疹/HPV/梅毒。
  \item 若有体液交换（如经血），注意屏障或事后清洁。
  \item \textbf{安全套可用于此类接触}（套住生殖器），能降低体液传播的风险。
\end{itemize}

\subsection*{共用器具（如有）}

\begin{itemize}[leftmargin=2em]
  \item 器具上残留的体液\textbf{可传播 STI}。
  \item 用后\textbf{清洁/消毒}；或在器具上套安全套，换人/换部位时\textbf{换新套}；或用一次性器具。
  \item CDC 明确：安全套在共用器具上"高效"预防 HIV 和部分 STI。
\end{itemize}

\section{通用防护清单}

\begin{enumerate}[leftmargin=2em,itemsep=0.3em]
  \item \textbf{保持清洁}：手、器具、身体。
  \item \textbf{避开活动性病变}：有疱疹水疱、HPV 疣、溃疡时，避免皮肤--皮肤密切接触。
  \item \textbf{不共用未清洁器具}；共用时在器具上套安全套并更换。
  \item \textbf{有可见血液/牙龈出血/口腔破损时}，避免深吻和口腔接触。
  \item \textbf{疫苗是基础}：HPV、乙肝疫苗。
  \item \textbf{双方检测是根本}：多数 STI 无症状，"看起来健康"说明不了什么。
\end{enumerate}

\section{沟通：偏好和边界照样要谈}

非插入类的偏好同样需要表达和协商：

\begin{itemize}[leftmargin=2em]
  \item 有人\textbf{只想要非插入}，有人想探索更多——都是正当的偏好。
  \item \textbf{主动说偏好，别让对方猜}："我喜欢被慢一点""这里我可以，那里我先不要"。
  \item 边界照样成立：任何不适、任何"不"，都立刻尊重。
\end{itemize}

\section{什么时候需要额外注意}

\begin{itemize}[leftmargin=2em]
  \item 出现\textbf{活动性疱疹水疱、HPV 疣、溃疡}——发作期传播风险显著升高，先避免接触，咨询医生。
  \item 有\textbf{可见血液、牙龈出血、口腔破损}时，避免深吻和口腔接触。
  \item \textbf{共用器具}且对方状态不明——按"清洁/套安全套/换套"处理，必要时双方检测。
  \item 若不确定是否感染，按第~\ref{chap:sti}~章安排检测。
\end{itemize}

\section{给小白的一页清单}

\begin{enumerate}[leftmargin=2em,itemsep=0.3em]
  \item 我把它当作\textbf{完整的亲密行为}，不是"前戏附属"。
  \item 我知道 HIV 主要看体液（非插入类风险极低），但\textbf{疱疹/HPV/梅毒看皮肤接触}——这才是重点。
  \item 接触前洗手、指甲干净，避免刮伤黏膜。
  \item 手在不同部位间移动时，洗手或戴一次性手套。
  \item 有活动性疱疹/疣/溃疡时，避免皮肤--皮肤密切接触。
  \item 不共用未清洁器具；共用时套安全套并更换。
  \item 有可见血/牙龈出血/口腔破损时，避免深吻和口腔接触。
  \item 我主动谈自己的偏好和边界，尊重对方的"不"。
  \item 疫苗（HPV/乙肝）+ 双方检测是根本。
\end{enumerate}

\chapter{非插入类行为的愉悦技巧：从"不危险"到"很愉悦"}
\label{chap:outer-pleasure}

前一章讲了这些行为\textbf{如何防病}。这一章换个角度：\textbf{这些行为本身就是完整的性体验，而且对很多人来说，比插入更容易放松、更接近愉悦。} 先把最关键的一个认知立住——

\begin{center}
\myhei\large 对多数女性而言，愉悦的核心不在"里面"，而在\textbf{阴蒂}。
\end{center}

数据支持：约 70--80\% 的女性需要\textbf{直接的阴蒂刺激}（手指/口/摩擦）才能达到高潮；而仅靠阴道插入、缺少阴蒂刺激时，高潮率仅约 21--30\%。也就是说——\textbf{非插入类行为不是"插入的替代品"，它往往就是"高潮的主通道"。}

所以"非插入"不等于"差一档"，它只是\textbf{刺激点不同}。下面逐类讲怎么做。

\section{贯穿所有技巧的底层逻辑：唤起是慢的}

\begin{itemize}[leftmargin=2em,itemsep=0.3em]
  \item 女性从平静到充分唤起，生理上常常需要 \textbf{20--45 分钟}（男性通常更快）。"心理唤起必须走在生理前面"——\textbf{大脑是最大的性器官}。
  \item 推论：所有技巧都服务于"把唤起的时间拉长、把紧张降下来"。\textbf{慢不是偷懒，慢本身就是技巧。}
\end{itemize}

\section{逐类技巧}

\subsection*{1. 全身爱抚（不只是"前戏"）}

\textbf{要点：从远到近、从慢到快、从大到小。}

\begin{itemize}[leftmargin=2em,itemsep=0.3em]
  \item \textbf{敏感区分布}（先避开"重点"）：耳后与耳垂、颈侧、锁骨与胸口、内臂内侧、下背与腰窝、大腿内侧、腰臀。这些区域神经末梢密、心理上也"安全"，适合开场。
  \item \textbf{力度节奏}：先轻（指尖/手掌大面积轻抚），观察反应，再逐渐加力、缩小范围。\textbf{力度的变化}比力度的大小更重要——忽快忽慢、忽轻忽重比"一直猛"更让人兴奋。
  \item \textbf{动作细节}：手掌平贴、大面积移动，让皮肤有"预期"；指尖点触、画圈，制造"意外"。别急着碰最敏感区——先让身体整体热起来。
  \item \textbf{怎么判断对方享受}：呼吸变深变快、身体不自觉贴近、轻哼、手主动搭上你。这些是"绿灯"信号。
\end{itemize}

\subsection*{2. 亲吻}

\textbf{不只是嘴唇，是一整套。}

\begin{itemize}[leftmargin=2em,itemsep=0.3em]
  \item \textbf{激活放松反应}：深而慢的吻会激活副交感神经，降低紧张与心跳焦虑——这是"把身体从战斗模式切换到休息模式"的开关，对容易紧张的人尤其重要。
  \item \textbf{技巧}：从额头、脸颊、嘴角开始，再过渡到唇，\textbf{别一上来就猛吻}；唇的轻碰→轻咬→舌的轻触，\textbf{由浅入深}，配合对方节奏。
  \item 颈、耳后、锁骨是吻的延伸区，很多人这里比嘴更敏感。
  \item \textbf{常见误区}：只顾自己"技术"，不看对方——\textbf{对方的呼吸与身体前倾是最好的节拍器}。
\end{itemize}

\subsection*{3. 手指刺激（对多数人是最直接的"非插入"高潮通道）}

\textbf{核心：慢 + 润滑 + 耐心，而不是快。}

\begin{itemize}[leftmargin=2em,itemsep=0.3em]
  \item \textbf{阴蒂是重点}：用\textbf{指腹}（不是指甲）在阴蒂头及周围\textbf{小范围、慢速度}打圈或上下轻抚。压力要轻——阴蒂极敏感，太用力反而发麻甚至疼。
  \item \textbf{润滑很重要}：天然润滑可能不够，\textbf{用润滑剂能显著提升舒适度}（水基最通用）。
  \item \textbf{节奏}：接近高潮时\textbf{放慢、减力}，让对方在高潮边缘停留（这叫 \emph{edging/边缘控制}），往往比"直线冲顶"更强烈、更持久。
  \item \textbf{前壁按压}："come hither"手势（食指中指轻弯向上），力度轻、节奏稳；有人喜欢有人没感觉，\textbf{问比猜准}。
  \item \textbf{男性}：手指包裹、刺激冠状沟与系带、龟头颈部，同样讲究慢和润滑。
  \item \textbf{卫生}：指甲修短剪平、手部清洁；进出肛门后务必更换或彻底清洗并换套（与前一章一致）。
\end{itemize}

\subsection*{4. 口部（对阴蒂 / 对生殖器）}

\textbf{要点：放松 + 观察 + 让对方主导节奏。}

\begin{itemize}[leftmargin=2em,itemsep=0.3em]
  \item \textbf{对阴蒂}：以\textbf{舌头轻抚为主}，小范围、慢速；唇轻碰、舌头打圈，\textbf{别用牙齿硬顶}。很多人喜欢"\textbf{轻触阴蒂 + 同时手指进入}"的\textbf{双重刺激}，比单一动作更接近高潮。
  \item \textbf{对男性}：\textbf{从下往上包裹、轻触系带与冠状沟}，力度要轻（龟头很敏感，太用力反而麻）；\textbf{节奏稳定}比快慢剧烈更好，接近高潮时放慢。
  \item \textbf{共同要点}：放松嘴部肌肉、用身体（不只用嘴）参与、\textbf{别憋气}（憋气=你自己先累了）。
  \item \textbf{安全}：口--口/口--生殖器可传播疱疹(HSV)、HPV、淋病、梅毒；有口腔/生殖器溃疡、疱疹活动时避免；\textbf{口腔也可用口用安全套或牙科薄膜(dental dam)}。
\end{itemize}

\subsection*{5. 摩擦 / 挤压（外--外、身体贴合）}

\textbf{这是最"安全"也最被低估的一类。}

\begin{itemize}[leftmargin=2em,itemsep=0.3em]
  \item \textbf{外--外（genital-to-genital）}：两人身体贴合、生殖器相互摩擦——\textbf{没有体液交换、几乎无感染风险，却能得到很强的刺激}。是新手、建立身体默契的最佳起点。
  \item \textbf{腿夹/体位贴合}：侧卧相对、双腿交叠，靠身体移动产生摩擦。
  \item \textbf{技巧}：\textbf{慢移}大于快速，用骨盆轻蹭而非大幅扭动；重点是"贴合面积与角度"，让对方找到最舒服的角度并停住。
\end{itemize}

\subsection*{6. 玩具辅助（可选，非必需）}

\begin{itemize}[leftmargin=2em,itemsep=0.3em]
  \item 振动器对阴蒂的刺激，对"手指够不到"的很多人是\textbf{关键}（持续振动频率是手指难以做到的）。
  \item \textbf{使用要点}：先用\textbf{最低档}、从外围靠近；振动直接压在阴蒂上容易"麻"，\textbf{隔着振动头轻触或稍偏一点}往往更舒服；事后清洗消毒。
\end{itemize}

\section{贯穿所有技巧的三件"元技巧"}

\begin{enumerate}[leftmargin=2em,itemsep=0.3em]
  \item \textbf{沟通是最高级的技巧}——一句"舒服吗""轻一点？""快一点还是慢一点？"比任何"技术"都有效。这不是打断情趣，而是\textbf{把命中率翻倍}。
  \item \textbf{读身体信号}（对方呼吸、贴靠、声音、手的反应）往往比问还准——但关键节点还是开口问。
  \item \textbf{把高潮当过程而非 KPI}（见\ref{chap:cycle}）：越放松、越不"赶着到终点"，体验越好。
\end{enumerate}

\section{安全要点（简短，不展开）}

\begin{itemize}[leftmargin=2em,itemsep=0.3em]
  \item 口--口 / 口--生殖器：疱疹、HPV、淋病、梅毒可经黏膜传播；\textbf{有溃疡/疱疹活动时避免}。
  \item 任何可能接触肛门的动作：前后必须换套/彻底清洗。
  \item 手部接触前清洁、指甲修平。
  \item 工具：水基润滑通用、事后清洗消毒。
\end{itemize}

\section{一页速查}

\begin{tabularx}{\textwidth}{lCC}
\toprule
行为 & 愉悦核心 & 一句话技巧 \\
\midrule
全身爱抚 & 全身敏感区 & 由远到近、力度变化 \\
亲吻 & 放松+全身 & 由浅入深、看对方 \\
手指 & 阴蒂/前壁 & 慢+润滑+边缘控制 \\
口部 & 阴蒂/系带 & 轻触为主、放松别憋气 \\
摩擦挤压 & 外--外贴合 & 慢移找角度 \\
玩具 & 阴蒂振动 & 低档起步、轻触 \\
\bottomrule
\end{tabularx}

\vspace{1em}
\noindent\textbf{一句话总结：} 非插入类行为的技巧精髓 = \textbf{慢 + 润滑 + 从外到内 + 全程读对方反应 + 关键处开口问}，核心刺激点是\textbf{阴蒂}而非"里面"。

\chapter{亲密关系沟通：从"会说话"到"更亲密"}
\label{chap:communication}

前面的章节反复出现同一个词——\textbf{沟通}。这一章把它单独展开：因为\textbf{沟通不是"技巧"里的一个选项，而是让所有技巧生效的开关}。很多尴尬、误判、"为什么我没感觉"，根子都在"没说"或"没听懂"。

\section{一个最好用的框架：非暴力沟通（NVC）四步}

美国心理学家 Marshall Rosenberg 提出的非暴力沟通，可浓缩为四步——\textbf{观察、感受、需要、请求}。关键是\textbf{用"事实+感受"代替"评价+指责"}。

\begin{tabularx}{\textwidth}{p{2.6cm}XX}
\toprule
步骤 & 怎么说（示例） & 避免（反面） \\
\midrule
观察（事实，不带评判） & "刚才那一下你好像顿了一下" & "你又在敷衍我" \\
感受（我的情绪） & "我有一点点失落" & "你根本不在乎" \\
需要（背后的诉求） & "我其实想要更专注的感觉" & "你总是这样" \\
请求（具体、可执行） & "下次你能慢一点、看着我说吗" & "你以后能不能有点心" \\
\bottomrule
\end{tabularx}

\vspace{0.5em}
\noindent 核心区别：\textbf{请求}是可被回答"是/否"的具体动作；\textbf{要求/指责}则会引发防御和反驳。

\section{四个高频场景的现成话术}

\begin{enumerate}[leftmargin=2em,itemsep=0.3em]
  \item \textbf{过程中}（微调，不打断节奏）："舒服吗？""轻一点 / 重一点？""再靠近一点。"——一句话就把命中率拉高。
  \item \textbf{设边界}（现在不想）："我暂时不想做这个，换一个可以吗？"——\textbf{说"不"不需要给理由}，也不代表否定对方这个人。
  \item \textbf{事前}（我想要什么）："我最想要的是……你愿意试试吗？"
  \item \textbf{事后}（回顾）："刚才哪部分你最喜欢？""有没有哪里其实不太舒服？"——把体验变成下次的进步。
\end{enumerate}

\section{倾听：比会说更重要}

\begin{itemize}[leftmargin=2em,itemsep=0.3em]
  \item \textbf{复述确认}："你的意思是……对吗？"——确认没听偏，也让人被听见。
  \item \textbf{不急于反驳或辩解}：先接住，再表达自己。
  \item \textbf{注意非语言信号}：对方犹豫、沉默、身体往后，可能比嘴上说的更真实——\textbf{有疑就问，别替对方下结论}。
\end{itemize}

\section{时机：什么时候聊最合适}

\begin{itemize}[leftmargin=2em,itemsep=0.3em]
  \item \textbf{事后放松时} > 过程中 > 争吵后。刚结束、情绪平稳时聊感受，接受度最高。
  \item \textbf{事前约定}比事后补救更省力（尤其避孕、检测、边界）。
  \item \textbf{避开}酒精影响判断、疲惫、正在生气这三个"低质量沟通"时段。
\end{itemize}

\section{五个最常见的沟通坑}

\begin{enumerate}[leftmargin=2em,itemsep=0.3em]
  \item 把"我不喜欢"说成"你错了"——从感受变成攻击。
  \item 用沉默代替表达，指望对方"自己懂"——但没有人是读心师。
  \item 事后从不回顾，体验无法迭代。
  \item 把对方的身体反应当"同意"，其实\textbf{持续、自愿、可撤回}才是同意的核心。
  \item 用"随便/都行"逃避选择——看似随和，实则把压力丢给对方。
\end{enumerate}

\section{三个能长期受益的习惯}

\begin{enumerate}[leftmargin=2em,itemsep=0.3em]
  \item \textbf{事前问}：先谈清楚想要什么、边界在哪。
  \item \textbf{过程中看反应}：把对方的身体当成最准的反馈器。
  \item \textbf{事后聊感受}：一句"哪部分你喜欢/不喜欢"胜过无数次猜测。
\end{enumerate}

\section{什么时候该寻求专业帮助}

如果沟通长期困难、存在语言/冷暴力，或涉及创伤、强迫、焦虑（如"一想到就紧张"），这已经不是"学点话术"能解决的——\textbf{找正规医疗机构的性医学科、心理科或正规心理咨询}，是专业且被鼓励的选择，不是"有病"。

\section{一页清单}

\begin{enumerate}[leftmargin=2em,itemsep=0.4em]
  \item 我用"事实+感受+请求"表达，而不是评价与指责。
  \item 我说"不"时不必给理由，也尊重对方的"不"。
  \item 我会在事后用一两句话回顾"喜欢/不喜欢"。
  \item 我把对方的身体反应当作最准的反馈器，有疑就问。
  \item 我在情绪平稳、清醒时谈重要话题，避开酒精与争执。
  \item 我接受"没人是读心师"——重要的事都要说出口。
\end{enumerate}

\chapter{避孕方法详解：怎么选才靠谱}

\begin{center}
\begin{tikzpicture}[font=\small]
  \node[anchor=east] at (0,0) {IUD / 埋植棒(LARC)};
  \fill[mygreen] (0.2,0) rectangle (0.247,0.3);
  \node[mygray] at (0.75,0.15) {\textless 1\%};
  \node[anchor=east] at (0,-0.7) {结扎};
  \fill[mygreen] (0.2,-0.7) rectangle (0.247,-0.4);
  \node[mygray] at (0.75,-0.55) {\textless 1\%};
  \node[anchor=east] at (0,-1.4) {口服药 / 贴片 / 环};
  \fill[myteal] (0.2,-1.4) rectangle (0.62,-1.1);
  \node[mygray] at (1.1,-1.25) {约 9\%};
  \node[anchor=east] at (0,-2.1) {安全套(男用)};
  \fill[myteal] (0.2,-2.1) rectangle (0.81,-1.8);
  \node[myteal] at (1.5,-1.95) {约 13\% · 唯一防 STI};
  \node[anchor=east] at (0,-2.8) {安全期 / 日历法};
  \fill[myamber] (0.2,-2.8) rectangle (1.05,-2.5);
  \node[mygray] at (1.55,-2.65) {约 18\%};
  \node[anchor=east] at (0,-3.5) {体外射精};
  \fill[myamber] (0.2,-3.5) rectangle (1.14,-3.2);
  \node[mygray] at (1.64,-3.35) {约 20\%};
  \node[anchor=east] at (0,-4.2) {安全套 + 杀精剂};
  \fill[myamber] (0.2,-4.2) rectangle (1.52,-3.9);
  \node[mygray] at (2.0,-4.05) {约 28\%};
  \node[anchor=east] at (0,-4.9) {不避孕};
  \fill[myred] (0.2,-4.9) rectangle (4.2,-4.6);
  \node[myred] at (4.7,-4.75) {约 85\%};
\end{tikzpicture}
\par\vspace{0.3em}\noindent\footnotesize 图：各避孕方法典型使用失败率（每年）。条形越长越易意外怀孕。
\end{center}

\label{chap:contraception}

前面多次提到安全套与紧急避孕。这一章把\textbf{全部避孕方法}放进同一张图里看，回答一个最实际的问题：\textbf{怎么选才既防怀孕、又防性病。}

\section{先建立一个关键认知：两种"失败率"差别巨大}

\begin{itemize}[leftmargin=2em,itemsep=0.3em]
  \item \textbf{完美使用}：每次都严格按说明做（每天按时吃、每次都戴）。
  \item \textbf{典型使用}：真实生活里——会漏服、会忘、会中途出问题。
\end{itemize}

\noindent 同一个方法，两个数字可能差 3 倍以上。判断"靠不靠谱"，要看\textbf{典型使用失败率}（按真实使用，一年里约有多少人意外怀孕）。

\section{按"真实可靠度"分层（越靠上越可靠）}

\begin{tabularx}{\textwidth}{p{3.4cm}CCC}
\toprule
方法 & 典型使用约怀孕比例(每年) & 能否防性病 & 是否长效 \\
\midrule
宫内节育器/皮下埋植(LARC) & $<$1\% & 否 & 是(多年) \\
男/女结扎 & $<$1\% & 否 & 是(不可逆) \\
短效口服药/贴片/阴道环 & 约 7--12\% & 否 & 否(需每天) \\
安全套(男用) & 约 13\% & \textbf{是(唯一日常)} & 否(每次) \\
安全套+杀精剂 & 约 28\% & 部分 & 否(每次) \\
安全期/日历法 & 约 12--24\% & 否 & 否(每次) \\
体外射精 & 约 20\% & 部分(前液) & 否(每次) \\
不避孕 & 约 85\% & — & — \\
\bottomrule
\end{tabularx}

\vspace{0.5em}
\noindent\textbf{一句话：} 要"几乎不担心怀孕"选 LARC（IUD/埋植）；要"顺便防性病"必用安全套。\textbf{两者不冲突，可以叠加}（安全套 + 口服药 = 既防孕又防 STI）。

\section{几个高频方法的要点}

\subsection*{1. 安全套（唯一日常防 STI 的方法）}

防 STI 能力取决于"每次都正确用"：\textbf{事前戴、全程戴、不翻转、戴在勃起上、射精后按住退出、用完即弃}。漏用/滑脱是中途出问题的主因；\textbf{不要与油基润滑同用}（会溶乳胶，见\ref{chap:practices}）。

\subsection*{2. 短效口服避孕药（需每天、防不了 STI）}

核心是"每天都吃、别漏"——\textbf{漏服/晚服是失败主因}。常见误区：刚开始吃的前 7 天要配合安全套；\textbf{它不防任何性病}。需医生评估适应证（偏头痛伴先兆、吸烟且年龄较大、血栓风险等人群多不宜）。

\subsection*{3. LARC（长效可逆：IUD / 埋植棒）——一次放置、多年有效}

放置后基本"免维护"，典型失败率 $<$1\%，是目前\textbf{非绝育里最可靠}的。\textbf{可逆}：取出后生育力通常很快恢复。需到正规医院放置/取出；"上环影响生育"多为误解——取出后不影响。

\subsection*{4. 紧急避孕（事后补救，不是日常方法）}

\begin{itemize}[leftmargin=2em,itemsep=0.3em]
  \item 左炔诺孕酮片：\textbf{越早越好，72 小时内}效果最佳（5 天内仍部分有效）。
  \item 乌利司坦：120 小时内。
  \item 铜 IUD：\textbf{事后 5 天内放置，是最可靠的紧急避孕}。
  \item \textbf{不能替代常规避孕}，频繁用会紊乱经期、失败率高于常规法。
\end{itemize}

\section{小白最容易踩的几个坑}

\begin{enumerate}[leftmargin=2em,itemsep=0.3em]
  \item \textbf{"第一次/月经前后是安全期"}——排卵会提前、精子可存活数天，安全期并不可靠。
  \item \textbf{体外射精很安全}——射精前已有含精子/STI 的前液，典型失败率约 20\%。
  \item \textbf{"吃了药就不用套了"}——口服药只防孕，不防 STI；性伴不明确时两者都要。
  \item \textbf{靠"感觉"判断排卵期}——身体感受不准，别赌。
\end{enumerate}

\section{怎么选（决策逻辑）}

\begin{itemize}[leftmargin=2em,itemsep=0.3em]
  \item \textbf{要不要防 STI？} → 要：\textbf{安全套是底线}。
  \item \textbf{想长期、几乎不想操心？} → LARC（IUD/埋植），需就医。
  \item \textbf{暂时、短期、可逆、能按时吃药？} → 短效口服药 + 安全套。
  \item \textbf{一次性、不确定长期关系？} → 安全套为主。
  \item \textbf{已发生无保护行为？} → 尽快紧急避孕（72h 内）+ 考虑检测。
\end{itemize}

\vspace{1em}
\noindent\textbf{一句话总结：} 最可靠的组合 = \textbf{LARC/口服药管怀孕 + 安全套管性病}；体外射精和安全期都不可靠；紧急避孕是补救不是日常。

\chapter{口交、阴道性交、肛交：逐类深入}

\begin{center}
\begin{tikzpicture}[
  cell/.style={draw=black!40,minimum width=3.0cm,minimum height=0.85cm,align=center,font=\small},
  head/.style={cell,font=\bfseries}
]
\matrix[
  matrix of nodes,
  nodes in empty cells,
  row sep=-0.4pt, column sep=-0.4pt
]{
  \node{}; & \node[head,fill=myblueL]{口交}; & \node[head,fill=mypurpleL]{阴道性交}; & \node[head,fill=mycoralL]{肛交}; \\
  \node[cell,align=left]{会怀孕？}; & \node[cell]{否}; & \node[cell,myred]{是}; & \node[cell]{否}; \\
  \node[cell,align=left]{STI 风险}; & \node[cell,fill=mygreenL]{中(咽部)}; & \node[cell,fill=myamberL]{高}; & \node[cell,fill=myredL]{最高}; \\
  \node[cell,align=left]{必用屏障}; & \node[cell]{dental dam\\口用套}; & \node[cell]{安全套}; & \node[cell]{安全套·强制}; \\
  \node[cell,align=left]{润滑}; & \node[cell]{少量即可}; & \node[cell]{需充分}; & \node[cell]{必须大量}; \\
  \node[cell,align=left]{关键铁律}; & \node[cell]{避免口--肛}; & \node[cell]{事后排尿防UTI}; & \node[cell]{前→后必换套}; \\
};
\end{tikzpicture}
\par\vspace{0.3em}\noindent\footnotesize 图：三类行为的风险与防护对照。
\end{center}

\label{chap:threeparts}

前面第\ref{chap:practices}章讲了这几类的健康框架。这一章逐类深入：每类都覆盖\textbf{怎么做更舒服 + 安全底线 + 常见误区}。定位仍是\textbf{医学科普 + 合理}——讲"怎么做、怎么安全"，不做露骨描写。

\section{一、口交（oral sex）}

\textbf{包含}：口对生殖器（男/女皆可）；以及\textbf{口--肛}（风险最高，单列）。

\textbf{怎么做更舒服}（与第\ref{chap:outer-pleasure}章一致）：以\textbf{舌头轻抚为主}，小范围、慢速，别用牙齿硬顶；\textbf{放松嘴部、用身体参与、别憋气}；全程看对方反应微调力度与节奏。

\textbf{风险（重点：口交不会怀孕，但会传 STI）}：
\begin{itemize}[leftmargin=2em,itemsep=0.3em]
  \item 口--生殖器可互染：\textbf{HSV-1（口唇疱疹→生殖器疱疹）、HPV、淋病、衣原体、梅毒}；咽部也可被淋病/衣原体感染且常无症状。
  \item \textbf{口--肛 = 粪口暴露}：可能传播大肠杆菌、志贺菌、溶阿米巴、甲肝等——\textbf{风险最高，建议避免}；若进行务必用套并事后更换。
\end{itemize}

\textbf{防护}：口腔可套\textbf{牙科薄膜(dental dam)或口用安全套}；\textbf{有活动口唇疱疹时避免对生殖器}；口交与肛交之间\textbf{必须换套/清洁}。

\textbf{常见误区}：「口交安全、不会得病」——不会怀孕，但\textbf{会传性病}，屏障工具同样适用。

\section{二、阴道性交（vaginal sex）}

\textbf{过程要点}：女性充分唤起常需 \textbf{20--40 分钟}，别急着进入——\textbf{先前戏、等自然润滑}；慢起步、观察对方，力度与节奏随反馈调整。

\textbf{避孕（最大风险点）}：\textbf{全程戴安全套}，或\textbf{安全套 + 口服药/LARC 叠加}；漏用、滑脱、事后才戴——这三样是意外怀孕与 STI 的主因。

\textbf{风险}：HIV、疱疹、HPV、淋病、衣原体、细菌性阴道病(BV) 等，是多数 STI 的主要传播途径。

\textbf{卫生 / 预防尿路感染(UTI)}：\textbf{事后尽快排尿 + 多饮水}（女性尿道短、易感染）；便后\textbf{由前向后擦}；前后清洁外生殖器即可，\textbf{不必冲洗阴道内部}（会破坏菌群）。

\textbf{常见误区}：「第一次不会怀孕/得病」「体内才会、体外就安全」——\textbf{怀孕与 STI 都看时机和防护，不看"第几次"}。

\section{三、肛交（anal sex）——风险最高，防护要求也最高}

\textbf{两个关键生理事实（决定一切）}：
\begin{enumerate}[leftmargin=2em,itemsep=0.3em]
  \item \textbf{直肠没有自然润滑}；
  \item \textbf{直肠黏膜薄、极易微裂}——微裂正是 HIV/STI 进入的通道。
\end{enumerate}

\textbf{防护（几乎是强制项）}：
\begin{itemize}[leftmargin=2em,itemsep=0.3em]
  \item \textbf{必须充分润滑}，且用\textbf{水基或硅基}（油基/凡士林会溶解乳胶套，使其失效）；
  \item \textbf{安全套必用}；
  \item \textbf{顺序铁律}：任何"前→后"（阴道→肛门、口→肛）\textbf{必须换新套并清洁}，避免把肛周细菌带到口或阴道；
  \item 从事肛交者建议接种\textbf{甲肝、乙肝疫苗}（CDC 推荐）。
\end{itemize}

\textbf{风险层级}：对 HIV/STI 而言，\textbf{接受方(receptive) 风险最高}，其次为插入方；同时传播疱疹、HPV、淋病、衣原体、甲肝。

\textbf{常见误区}：「洗干净就安全」——干净挡不住"黏膜薄 + 无自润滑"，仍属高风险；「肛交不伤身体」——反复微裂会持续增加感染风险。

\section{一页对照}

\begin{tabularx}{\textwidth}{lCCCC}
\toprule
项目 & 口交 & 阴道性交 & 肛交 \\
\midrule
会怀孕吗 & 否 & 是 & 否 \\
主要 STI 风险 & 中(含咽部) & 高 & 最高 \\
必用屏障 & dental dam/口用套 & 安全套 & 安全套(强制) \\
润滑 & 少量即可 & 需充分 & 必须大量(水/硅基) \\
关键铁律 & 避免口--肛 & 事后排尿防 UTI & 前→后必须换套 \\
\bottomrule
\end{tabularx}

\vspace{1em}
\noindent\textbf{三类的共同底线：} 屏障工具是\textbf{共同防线}，不是"可选"；\textbf{双方检测 + HPV/乙肝疫苗}才是根本；\textbf{沟通}（想要什么、边界在哪、是否检测过）是前提。

\noindent\textbf{一句话总结：} 口交防"黏膜互染"、阴道性交防"怀孕+STI"、肛交因"无自润滑+黏膜薄"而风险最高——\textbf{润滑足、安全套全程、前→后必换套、检测+疫苗兜底}，这几类行为才能真正既愉悦又安全。

\chapter{阴道性交详解（单独深入）}
\label{chap:vaginal}

\noindent\textbf{前提先立住：} 阴道性交的核心不是"深"，而是\textbf{唤起 + 阴蒂 + 放松}。多数女性的愉悦与高潮来自阴蒂而非"里面"（见\ref{chap:outer-pleasure}），所以"进去"只是过程的一部分，不是终点。

\section{生理基础：别急着"进入"}

\begin{itemize}[leftmargin=2em,itemsep=0.3em]
  \item 唤起时阴道会\textbf{伸展、湿润、充血}（医学称"tenting"，内部空间增大）——\textbf{需要一个过程}，先前戏、等它准备好，比硬塞更舒服也更安全。
  \item \textbf{干涩很常见}，原因：唤起不足、激素波动（经期/哺乳/更年期）、压力、某些药物（如抗组胺药）。\textbf{处理：用润滑剂}（水基，可配安全套）。
\end{itemize}

\section{完整过程（呼应\ref{chap:cycle}性反应周期）}

\begin{enumerate}[leftmargin=2em,itemsep=0.3em]
  \item \textbf{唤起期}：亲吻、爱抚、全身触摸——目标是充分湿润与放松，别跳过。
  \item \textbf{进入期}：慢、润滑足够、观察对方。
  \item \textbf{进行期}：节奏与角度随反馈调整。
  \item \textbf{高潮（非必须）}：多数女性需叠加阴蒂刺激。
  \item \textbf{消退期}：放松、事后护理。
\end{enumerate}

\section{进入的细节}

\begin{itemize}[leftmargin=2em,itemsep=0.3em]
  \item 角度与深度可随体位调整，\textbf{找到双方都舒服的位置停住}，比"换花样"更重要。
  \item \textbf{疼痛是信号，不是要克服的障碍}：疼就停下、回退、补润滑或先做前戏。\textbf{持续或剧痛要就医}（可能感染、阴道痉挛、激素问题），别硬忍。
  \item 润滑随需补，别等干涩才加。
\end{itemize}

\section{体位的"功能维度"（看需求，不追花样）}

用\textbf{谁主导 / 深度 / 节奏 / 刺激点}来理解，没有"最好"，只有"此刻适合什么"：

\begin{tabularx}{\textwidth}{lCCC}
\toprule
体位 & 特点 & 主导 & 适合 \\
\midrule
传教士(面对面) & 亲密、眼神交流 & 女方可控深度 & 想亲近、想沟通 \\
女上位 & 主导节奏/角度/深度、易兼顾阴蒂刺激 & 女方 & 想掌控、想兼顾高潮 \\
侧卧(勺子式) & 省力、温柔、压力最小 & 双方轻松 & 初次、累了、想放松 \\
后入 & 贴合较深、刺激点不同 & 男方多 & 想要更深贴合 \\
\bottomrule
\end{tabularx}

\vspace{0.5em}
\noindent 新手建议从\textbf{侧卧或传教士}起步——压力小、易放松、易沟通。

\section{节奏与力度}

\begin{itemize}[leftmargin=2em,itemsep=0.3em]
  \item \textbf{变化大于单一}：快慢、深浅交替，比一直猛更有感觉。
  \item \textbf{edging（边缘控制）}：接近高潮时放慢、退回，再推上去，往往更强烈持久。
  \item \textbf{叠加阴蒂刺激}（手/振动）——这是"只插不高潮"最直接的解法。
\end{itemize}

\section{避孕（必做，最大风险点）}

\begin{itemize}[leftmargin=2em,itemsep=0.3em]
  \item \textbf{全程安全套}，或\textbf{安全套 + 口服药/LARC 叠加}。
  \item 漏用、滑脱、事后才戴 = 意外怀孕与 STI 的三大主因。
  \item 安全套是唯一既防孕又防 STI 的日常方法（见\ref{chap:contraception}）。
\end{itemize}

\section{卫生 / 预防尿路感染(UTI)}

\begin{itemize}[leftmargin=2em,itemsep=0.3em]
  \item \textbf{事后尽快排尿 + 多饮水}（女性尿道短，性行为易把细菌带入）。
  \item 便后\textbf{由前向后擦}。
  \item 前后清洁\textbf{外生殖器}即可，\textbf{不要冲洗阴道内部}（破坏菌群，反而易感染）。
\end{itemize}

\section{几个高频困惑（很多是"正常"，别自我怀疑）}

\begin{enumerate}[leftmargin=2em,itemsep=0.3em]
  \item \textbf{第一次会痛吗？} 轻度不适常见（紧张、润滑不足）；\textbf{剧痛不正常}（未充分唤起、阴道痉挛、感染）——别忍，处理或就医。
  \item \textbf{高潮不来是不是我有问题？} 多数女性需直接阴蒂刺激才到高潮，"光插不够"非常普遍，\textbf{不代表异常}。
  \item \textbf{"深才舒服"吗？} 这是误区——很多人更喜欢\textbf{浅而密 + 阴蒂刺激}，深不等于舒服。
  \item \textbf{干涩 = 没感觉？} 不一定，干涩有激素/药物/压力等多种原因，\textbf{润滑剂是正常工具，不是"不行"的标志}。
\end{enumerate}

\section{沟通（命中率翻倍的关键）}

关键节点一句话就够：\textbf{"深度合适吗？""快一点 / 慢一点？""疼不疼？""舒服吗？"} 把对方的呼吸、贴靠、声音当作最准的反馈器。

\vspace{1em}
\noindent\textbf{一句话总结：} 阴道性交 = \textbf{唤起到位 + 润滑足够 + 慢起步找角度 + 叠加阴蒂刺激 + 全程避孕 + 事后防 UTI}；"深"和"高潮"都不是目标，\textbf{放松、愉悦、安全}才是。

\chapter{常用性交姿势：按功能选}
\label{chap:positions}

\noindent\textbf{核心观念：} 没有"最好的姿势"，只有\textbf{此刻适合什么}；姿势是工具，服务于\textbf{唤起、舒适、愉悦和安全}，不追花样。

\section{按需求速查}

\begin{tabularx}{\textwidth}{C C}
\toprule
你想要 & 更合适的姿势 \\
\midrule
亲密、眼神交流 & 传教士位 \\
女方主导、兼顾高潮 & 女上位（正向/反向） \\
省力、轻柔、初次 & 侧卧位（勺子位） \\
更深贴合 & 后入式 \\
灵活、空间小 & 站立 / 坐姿 \\
孕期/疲劳/敏感 & 侧卧位、女上位 \\
\bottomrule
\end{tabularx}

\section{常用姿势逐个看}

\subsection*{1. 传教士位（男上、面对面）}
\begin{itemize}[leftmargin=2em,itemsep=0.3em]
  \item \textbf{功能}：面对面、便于眼神与肢体接触，\textbf{亲密感最强}；男方较省力、可主导深度节奏；女方腿搭肩/腰或垫高臀部可加深。
  \item \textbf{兼顾阴蒂}：面对面，\textbf{双方都方便用手或振动同时刺激阴蒂}——对"想兼顾高潮"的人很友好。
  \item \textbf{适合}：初次、想亲近、想交流、体力一般。
  \item \textbf{注意}：对男方腰腹负担较大；女方腰部不适可在臀下垫枕头。
\end{itemize}

\subsection*{2. 女上位（woman on top）}

分\textbf{正向}（面对面）与\textbf{反向}（背对）。
\begin{itemize}[leftmargin=2em,itemsep=0.3em]
  \item \textbf{功能}：\textbf{女方完全主导}深度、角度、节奏——最容易自我调节、也最易兼顾阴蒂角度。
  \item \textbf{适合}：想女方掌控、想兼顾高潮、男方想省力、想清楚表达"快/慢/深/浅"。
  \item \textbf{注意}：需要一定稳定性与体力；起身、落座时留点余地。
\end{itemize}

\subsection*{3. 侧卧位 / 勺子位（side-by-side）}
\begin{itemize}[leftmargin=2em,itemsep=0.3em]
  \item \textbf{功能}：\textbf{最省力、最温柔、压力最小}；可从背后贴合，深度较浅、撞击轻。
  \item \textbf{适合}：初次、敏感、疲劳、想要轻柔、\textbf{孕期较舒适}；也适合"聊着天慢慢来"。
  \item \textbf{注意}：深度浅，想要更深可配合手/振动刺激。
\end{itemize}

\subsection*{4. 后入式（后位）}
\begin{itemize}[leftmargin=2em,itemsep=0.3em]
  \item \textbf{功能}：贴合较深、角度不同、刺激点不同；节奏可快可慢。
  \item \textbf{适合}：想要更深贴合、想换一种刺激点。
  \item \textbf{注意}：\textbf{深浅要沟通}，过深过猛会撞到宫颈/造成不适；新手别一上来就猛。
\end{itemize}

\subsection*{5. 站立 / 坐姿}
\begin{itemize}[leftmargin=2em,itemsep=0.3em]
  \item \textbf{功能}：灵活、随时、空间小也能做；可借助墙、桌沿、椅子。
  \item \textbf{适合}：想快、想换环境、体力好的时候。
  \item \textbf{注意}：支撑与安全（别滑倒）、女方久站易累，可坐椅子。
\end{itemize}

\subsection*{6. 骑跨变体（女下、双腿交叠或并拢）}
\begin{itemize}[leftmargin=2em,itemsep=0.3em]
  \item \textbf{功能}：贴合深、便于控制角度；并拢腿偏浅，交叠/搭肩偏深。
  \item \textbf{适合}：想在"女上"基础上让男方主导一点、又想贴合。
\end{itemize}

\section{选姿势的 4 个关键维度}

\begin{enumerate}[leftmargin=2em,itemsep=0.3em]
  \item \textbf{谁主导}：想清楚"此刻谁来决定节奏"，会大幅减少尴尬。
  \item \textbf{深还是浅}：深（后入、双腿交叠）刺激点不同；浅（侧卧、女上轻压）更温柔。
  \item \textbf{要不要兼顾阴蒂}：面对面（传教士、女上正向）最方便叠加阴蒂刺激。
  \item \textbf{体力/疲劳/敏感度}：累、敏感、初次 → 侧卧、女上；体力好、想贴合 → 后入、骑跨。
\end{enumerate}

\section{健康提醒（任何姿势都适用）}

\begin{itemize}[leftmargin=2em,itemsep=0.3em]
  \item \textbf{润滑}：唤起不足就干涩，\textbf{用润滑剂}（水基可配安全套）——别等干涩才加。
  \item \textbf{避孕}：任何姿势都要\textbf{全程安全套}或叠加口服/LARC，姿势不改变怀孕与 STI 风险。
  \item \textbf{疼痛 = 停}：不适就停下、回退、补润滑、改姿势；持续剧痛要就医（感染/痉挛/激素等）。
  \item \textbf{沟通}：一句"深浅合适吗？""快一点/慢一点？""疼不疼？"比换十个姿势都管用。
\end{itemize}

\vspace{1em}
\noindent\textbf{一句话总结：} 姿势不是"比谁花样多"，而是\textbf{按当下需求（主导、深浅、阴蒂兼顾、体力）选工具}；配合润滑、避孕、沟通，任何姿势都能既舒服又安全。

\chapter{特殊场景：孕期与初次的渐进方案}

这两个场景的共同点是：\textbf{身体和心理都需要更慢、更温柔、更多沟通}。上一章讲的是"一般情况怎么选姿势"，本章讲"特殊情况下怎么调整"。

\section{孕期：安全地保持亲密}

\subsection{核心认知（先安心）}

\begin{itemize}[leftmargin=2em,itemsep=0.3em]
  \item \textbf{除非医生/助产士明确说不要，孕期性生活是安全的}——宝宝在子宫里有羊水和子宫壁保护，不会伤到（NHS/Tommy's）。
  \item 孕期性欲变多或变少\textbf{都正常}；亲吻、按摩、相互爱抚同样算亲密，不是"只有插入才算"。
  \item 孕晚期性交或高潮后可能出现 \textbf{Braxton Hicks（假宫缩，肚子发硬）}——这是正常现象，\textbf{不是临产}。
\end{itemize}

\subsection{哪些情况需要避免（医学红线）}

\begin{tabularx}{\linewidth}{X C}
\toprule
\textbf{情况} & \textbf{原因} \\
\midrule
前置胎盘 / 低位胎盘 & 插入可能引发出血 \\
\textbf{胎膜已破（"破水"）} & 感染风险上升，\textbf{必须避免} \\
多次流产史 & 遵医嘱 \\
\textbf{任何不明原因阴道出血} & 先检查，检查前避免 \\
医生明确告知的其他情况 & 个体化处理，务必问清 \\
\bottomrule
\end{tabularx}

\vspace{0.5em}
\noindent 有 STI 症状或怀疑感染时：用安全套（包括口交）并去检测；不要与可能有 STI 的人发生关系。

\subsection{姿势怎么选}

\begin{itemize}[leftmargin=2em,itemsep=0.3em]
  \item \textbf{首选侧卧位（side-by-side / 勺子位）}：不压肚子、省力、温柔，孕中晚期最舒服；NHS 明确推荐"并排侧卧"作为尝试方向。
  \item \textbf{女上位}：女方掌控深度和节奏、肚子不受压、随时能停。
  \item \textbf{慎用男上女下（传教士）}：早期可能因乳房敏感不适，后期肚子碍事；若用，要偏一侧、用枕头垫住腹部、避免直接压迫。
  \item \textbf{深度与力度}：孕晚期宫颈更敏感，\textbf{浅一点、慢一点}，太深可能触发不适或宫缩。
  \item \textbf{支撑}：枕头垫在腹部下方或两腿之间，能明显省力。
  \item \textbf{别长时间仰卧压腹}：孕期本身建议 28 周后避免长时间仰卧（与血流/围产风险相关），亲密时也别长时间仰压。
\end{itemize}

\subsection{干涩与替代方式}

\begin{itemize}[leftmargin=2em,itemsep=0.3em]
  \item 激素变化常导致干涩——\textbf{用水基润滑剂}，别硬来。
  \item 医生建议避免插入时：\textbf{亲吻、拥抱、按摩、相互爱抚、玩具}都是选项（玩具要清洁，共用要换套，先问医生）。
  \item \textbf{口交一般安全，但任一方有口唇疱疹（cold sore）时要避免对生殖器}——同一种病毒，传给宝宝可能很严重。
\end{itemize}

\subsection{关于"性交会引产吗"}

\textbf{没有证据表明性交会帮助发动分娩，部分研究反而显示不会。} 想靠它催产不靠谱；但破水后确实必须避免。

\section{初次：把"第一次"变成渐进，而不是硬闯}

\subsection{核心认知（破除三个迷思）}

\begin{enumerate}[leftmargin=2em,itemsep=0.3em]
  \item \textbf{"处女膜"不是一层封闭的膜}——它是阴道口的黏膜褶皱（阴道瓣），中间本来就有孔，形态因人而异（环状、半月状、筛状等）。
  \item \textbf{初次不一定出血、也不一定痛}——很多女性完全不出血。"见红"不是"证明"，不出血也不代表什么。
  \item \textbf{疼痛和出血多来自可预防的原因}：紧张$\rightarrow$盆底紧绷、唤起不足$\rightarrow$干涩、动作太快太深$\rightarrow$黏膜小裂伤。
\end{enumerate}

\subsection{渐进五步（重点）}

\begin{tabularx}{\linewidth}{C X X}
\toprule
\textbf{步骤} & \textbf{做什么} & \textbf{要点} \\
\midrule
\textbf{0. 准备} & 私密舒适的环境；先说好"可以随时停"；备好安全套+水基润滑；别喝酒 & 酒精麻痹感觉、影响判断 \\
\textbf{1. 非插入亲密} & 拥抱、亲吻、全身爱抚 & 目标是放松和唤起，\textbf{不是赶进度} \\
\textbf{2. 等身体准备好} & 唤起常需 20--45 分钟；等自然润滑，不够就加润滑剂；手指先轻试慢适应 & 慢就是技巧 \\
\textbf{3. 进入} & 慢、浅、润滑足；选\textbf{侧卧}或\textbf{女上位} & 疼就停，别硬忍 \\
\textbf{4. 进行} & 节奏慢、幅度小；随时问"疼不疼/深不深/要不要慢点" & 不必追求高潮 \\
\textbf{5. 事后} & 尽快排尿防 UTI、清洁外阴、观察身体反应 & 呼应前面章节 \\
\bottomrule
\end{tabularx}

\subsection{正常 vs 需就医}

\begin{itemize}[leftmargin=2em,itemsep=0.3em]
  \item \textbf{轻微不适 / 少量点滴出血}：常见，观察即可。
  \item \textbf{剧痛、持续痛、大量出血、头晕}：就医（可能是裂伤等）。
  \item \textbf{反复尝试都剧痛}：可能是阴道痉挛或前庭痛——找妇科/性健康门诊，\textbf{不是"矫情"}，是可治疗的。
\end{itemize}

\subsection{常见误区}

\begin{itemize}[leftmargin=2em,itemsep=0.3em]
  \item "第一次一定会痛/出血" $\rightarrow$ 不对。
  \item "不出血就不是处女" $\rightarrow$ 错误观念，形态和弹性因人而异。
  \item "忍一忍就过去了" $\rightarrow$ 痛是信号，硬忍会加重恐惧和痉挛，形成恶性循环。
  \item "要一次到位" $\rightarrow$ 完全可以分几次慢慢来，亲密不是任务。
\end{itemize}

\subsection{安全底线（与全书一致）}

\textbf{全程安全套}（避孕+防 STI）、\textbf{沟通是核心技巧}、\textbf{疼痛=停}——这三条在任何阶段都不变。

\vspace{1em}
\noindent\textbf{一句话总结：} 孕期 = \textbf{先确认无医学禁忌 + 侧卧/女上位 + 润滑 + 慢而浅 + 破水必停}；初次 = \textbf{破除"必痛必出血"迷思 + 非插入起步 + 等唤起和润滑 + 侧卧/女上位 + 疼就停}。两者都不是"硬闯"，而是"慢下来、随时沟通"。

\chapter{生命周期：不同年龄段的性健康}

前面各章讲的是"通用原则"，本章补上唯一缺的维度——\textbf{身体会随年龄变化，性和亲密需求也要跟着调整}。分四个阶段：青春期、产后、更年期、老年。核心观念先立住：\textbf{性健康不是某一个年龄的事，而是贯穿一生的；每个阶段都有该关注的不同重点。}

\section{青春期：建立认知和保护的起点}

\textbf{身体变化}：第二性征出现、月经初潮/遗精、性欲萌动——全是正常生理，不必羞耻或恐慌。

\textbf{这个阶段最该建立的三个认知}：

\begin{enumerate}[leftmargin=2em,itemsep=0.3em]
  \item \textbf{避孕与防病是两件事}——WHO 明确：\textbf{安全套是唯一既能避孕、又能防 STI（含 HIV）的方法}；口服药、埋植、IUD 都只避孕不防病。
  \item \textbf{"第一次"没有标准剧本}——不出血、不痛都很常见（呼应前面章节）；别被成人片或同伴话术带节奏。
  \item \textbf{避孕不会导致不孕}——WHO 澄清：现代避孕方法（含 IUD、埋植）\textbf{不会造成不孕}；"上环影响生育"是常见误解。
\end{enumerate}

\begin{itemize}[leftmargin=2em,itemsep=0.3em]
  \item 越早建立"全程安全套 + 同意 + 检测"的习惯，越受益；青少年同样可安全使用避孕药和 LARC（需医生评估）。
  \item 遇到困惑找\textbf{正规渠道}（校医、疾控、正规医院妇科/皮肤科），别靠网络偏方。
  \item 遇到任何让你不舒服的触碰——\textbf{同意可以随时撤回}，这不是"半途而废"。
\end{itemize}

\section{产后：恢复到重新亲密}

\textbf{什么时候可以恢复同房}（NHS/WHO 共识）：

\begin{itemize}[leftmargin=2em,itemsep=0.3em]
  \item 一般建议\textbf{等到恶露停止、会阴/剖宫产伤口愈合}，通常是\textbf{产后约 6 周}；顺产有撕裂或剖宫产者按医生具体指导。
  \item 过早（恶露未净时）恢复会\textbf{增加感染风险}。
  \item 恢复初期\textbf{动作宜慢、宜浅、用润滑}（激素变化常致干涩，尤其哺乳期），疼就停。
\end{itemize}

\textbf{避孕——这是最容易被忽略的点}：

\begin{tabularx}{\linewidth}{X X}
\toprule
\textbf{要点} & \textbf{说明} \\
\midrule
\textbf{产后可能很快恢复排卵} & 月经没来 $\ne$ 不会怀孕；先排卵后月经，可能"没来月经就怀上" \\
\textbf{哺乳闭经避孕法（LAM）} & 需\textbf{同时满足}三条件：(1) 纯母乳喂养 (2) 月经尚未恢复 (3) 产后 6 个月内；任一条件改变即失效，失败率随之上升 \\
\textbf{药物避孕时机} & WHO：纯母乳喂养者避孕药一般\textbf{产后 6 个月}开始；部分母乳喂养者\textbf{产后 6 周}；\textbf{仅孕激素迷你 pill 产后可立即用、不影响乳汁} \\
\textbf{LARC（IUD/埋植）} & 高效、可逆，产后由医生评估放置时机 \\
\bottomrule
\end{tabularx}

\vspace{0.5em}
\noindent 一句话：\textbf{别把"哺乳期"当天然避孕}——要么严格满足 LAM 三条件，要么用可靠的避孕方法。

\section{更年期（围绝经期）：变化是正常的，可处理}

\textbf{身体变化}（通常 45--55 岁，因雌激素下降）：

\begin{itemize}[leftmargin=2em,itemsep=0.3em]
  \item \textbf{阴道干涩、灼热、同房痛}——称\textbf{泌尿生殖综合征（GSM）}，是雌激素下降的直接结果，\textbf{不是"你不想了"或"感情问题"}。
  \item 性欲可能波动；潮热、情绪、睡眠也会受影响。
  \item 这一变化\textbf{很常见、且可处理}，不必默默忍受。
\end{itemize}

\textbf{怎么处理（NHS 明确）}：

\begin{tabularx}{\linewidth}{X X}
\toprule
\textbf{方法} & \textbf{要点} \\
\midrule
\textbf{润滑剂 / 阴道保湿剂} & 第一步、最安全；同房用润滑，日常用保湿剂 \\
\textbf{阴道局部雌激素（local HRT）} & 低剂量、只作用于阴道，\textbf{不增加乳腺癌风险、通常无需配孕激素、可按需长期使用}；最多约 3 个月起效；需医生评估 \\
\textbf{全身 HRT} & 针对潮热/情绪/睡眠等，需医生权衡利弊；局部阴道症状优先用局部雌激素 \\
\textbf{睾酮（低性欲）} & 非更年期常规适应证，仅专科医生在特定情况（绝经后、HRT 无效的低性欲）考虑 \\
\bottomrule
\end{tabularx}

\vspace{0.5em}
\textbf{两个易错点}：

\begin{itemize}[leftmargin=2em,itemsep=0.3em]
  \item \textbf{局部雌激素 $\ne$ 全身 HRT}：前者风险低得多，别因为害怕 HRT 而连局部治疗也拒绝。
  \item \textbf{更年期 $\ne$ 不用避孕}：NHS 建议——\textbf{50 岁以下}停经后仍需避孕 2 年，\textbf{50 岁以上}需 1 年；\textbf{55 岁后}自然怀孕极罕见（仍需注意 STI 防护）。
\end{itemize}

\section{老年阶段：需求不消失，方式可调整}

\textbf{核心认知}：性需求不会因年龄"消失"——\textbf{亲密、触摸、陪伴的需要贯穿一生}。"老年人不该有性生活"既不科学也不人道。

\begin{itemize}[leftmargin=2em,itemsep=0.3em]
  \item \textbf{干涩更明显}：润滑剂和局部雌激素依然是首选（同上）。
  \item \textbf{男性 ED 更常见}：常与血管/糖尿病/药物相关，\textbf{是可治疗的}；PDE5 抑制剂需医生评估（尤其心血管用药者）。
  \item \textbf{慢性病与药物}：很多药（降压药、抗抑郁药等）会影响性功能——\textbf{可与医生沟通调整}，别自行停药。
  \item \textbf{安全套别停}：老年人 STI 感染率在部分人群上升，\textbf{新伴侣仍需防护}。
  \item \textbf{亲密形式可以更广}：拥抱、按摩、陪伴、相互爱抚——\textbf{"性"不等于"插入"}（呼应\ref{chap:outer}）。
\end{itemize}

\textbf{何时就医}：突发性功能变化、性交疼痛、出血、或因药物明显影响生活——都值得找医生，\textbf{不是"难为情"的事}。

\section{一页速览}

\begin{tabularx}{\linewidth}{C X X}
\toprule
\textbf{阶段} & \textbf{身体关键变化} & \textbf{最该记住的一件事} \\
\midrule
\textbf{青春期} & 第二性征、初潮/遗精、性欲萌动 & 安全套是唯一既避孕又防病；避孕不致不孕 \\
\textbf{产后} & 恶露、伤口、激素波动、干涩 & 月经没来 $\ne$ 不排卵；LAM 三条件缺一不可 \\
\textbf{更年期} & 雌激素降、GSM 干涩、性欲波动 & 局部雌激素安全有效；停经后仍需避孕一段时间 \\
\textbf{老年} & 干涩加重、ED、药物影响 & 需求正常、可治疗；新伴侣仍要防护 \\
\bottomrule
\end{tabularx}

\vspace{1em}
\noindent\textbf{一句话总结：} 生命周期视角 = \textbf{每个阶段身体不同、需求正常、都有对应办法}——青春期重"建立保护习惯"，产后重"恢复 + 别误判避孕"，更年期重"干涩可治、别怕局部雌激素"，老年重"需求不消失、可调整可治疗"。

\chapter{口交详解（单独深入）}
\label{chap:oral}

\noindent\textbf{定位：} 口交\textbf{不会怀孕，但会传 STI}——这是它和"安全"之间唯一的矛盾点，也是本章要讲清的重点。

\section{类型}

\begin{itemize}[leftmargin=2em,itemsep=0.3em]
  \item 口对男性生殖器、口对女性生殖器（皆属黏膜互染场景）。
  \item \textbf{口--肛}：风险最高的一类，\textbf{单独对待}（粪口暴露）。
\end{itemize}

\section{风险：别小看"咽部"和"口腔"这两个常被忽略的入口}

\begin{tabularx}{\textwidth}{lC X}
\toprule
途径 & 可传播 & 关键点 \\
\midrule
口--生殖器 & HSV-1、HPV、淋病、衣原体、梅毒 & 咽部可被淋病/衣原体感染，\textbf{常无症状} \\
口--肛 & 大肠杆菌、志贺菌、溶阿米巴、甲肝 & \textbf{粪口暴露}，细菌/寄生虫/病毒均可能 \\
\bottomrule
\end{tabularx}

\vspace{0.5em}
\begin{itemize}[leftmargin=2em,itemsep=0.3em]
  \item \textbf{HSV-1（口唇疱疹）→ 生殖器疱疹}：是最常见的"由口传到下身"的情形，即使没有明显疱疹也可能排毒。
  \item \textbf{HPV}：口腔/咽部也可感染，与口咽癌存在关联；\textbf{HPV 疫苗对男女都建议}。
  \item \textbf{梅毒}：口腔黏膜斑也是传染源。
  \item \textbf{HIV 经口传播极低}：除非口内/生殖器有可见出血 + 黏膜破损（见\ref{chap:sti}），一般不构成风险，无需为口交常规 PEP。
\end{itemize}

\section{防护（屏障工具同样适用）}

\begin{itemize}[leftmargin=2em,itemsep=0.3em]
  \item 对男性生殖器：可用\textbf{口用安全套}。
  \item 对女性生殖器：可用\textbf{牙科薄膜(dental dam)}（覆盖在生殖器上，口接触薄膜）。
  \item \textbf{有活动口唇疱疹时，避免对生殖器}。
  \item 口交与肛交之间\textbf{必须换套/清洁}，避免交叉。
  \item 多性伴或不确定对方感染状态时，屏障 + 定期检测更稳。
\end{itemize}

\section{怎么做更舒服（与第\ref{chap:outer-pleasure}章一致）}

以\textbf{舌头轻抚为主}、小范围慢速、别用牙齿硬顶；\textbf{放松嘴部、用身体参与、别憋气}；全程看对方反应微调力度与节奏；接近高潮时可放慢。

\section{误区与红旗}

\begin{enumerate}[leftmargin=2em,itemsep=0.3em]
  \item 「口交安全、不会得病」——不会怀孕，但\textbf{会传 STI}。
  \item 「没症状就没事」——咽部淋病/衣原体常无症状，需靠\textbf{检测}发现。
  \item \textbf{红旗}：咽部反复红肿疼痛、口腔/生殖器出现溃疡或水疱、排尿刺痛——尽早就医筛查。
\end{enumerate}

\vspace{1em}
\noindent\textbf{一句话总结：} 口交的关键不是"会不会怀孕"，而是\textbf{屏障工具（dental dam/口用套）+ 有疱疹时回避 + 定期检测}；咽部和口腔是常被漏掉的感染入口。

\chapter{肛交详解（单独深入）}
\label{chap:anal}

\noindent\textbf{核心生理（决定一切）：} \textbf{直肠没有自然润滑 + 黏膜薄、极易微裂}。微裂正是 HIV/STI 进入的通道——所以肛交的防护要求是所有行为里最高的。

\section{风险}

\begin{itemize}[leftmargin=2em,itemsep=0.3em]
  \item 对 HIV/STI：\textbf{接受方(receptive) 风险最高}，插入方次之。
  \item 可传播：HIV、疱疹、HPV、淋病、衣原体、甲肝、乙肝。
  \item 微裂（肛裂）反复存在时，感染风险进一步升高。
\end{itemize}

\section{防护铁律（几乎是强制项）}

\begin{enumerate}[leftmargin=2em,itemsep=0.3em]
  \item \textbf{必须充分润滑}，且用\textbf{水基或硅基}——\textbf{油基/凡士林会溶解乳胶套，使其失效}（见\ref{chap:contraception}）。
  \item \textbf{安全套必用}，且全程不破不脱。
  \item \textbf{顺序铁律}：任何"前→后"（阴道→肛门、口→肛）\textbf{必须换新套并清洁}。
  \item 从事肛交者\textbf{建议接种甲肝、乙肝疫苗}（CDC 推荐）。
  \item 多性伴或高危者，可与医生讨论 \textbf{PrEP}（HIV 暴露前预防用药）。
\end{enumerate}

\section{卫生}

\begin{itemize}[leftmargin=2em,itemsep=0.3em]
  \item 进入前可用温水/淋浴清洁肛周，\textbf{灌洗非必须}且过度反而刺激黏膜。
  \item 事后清洁肛周即可；\textbf{避免肛→口}，若进行务必用套并更换。
  \item 润滑随需补，别等痛了才加。
\end{itemize}

\section{疼痛与微裂（别硬忍）}

\begin{itemize}[leftmargin=2em,itemsep=0.3em]
  \item 轻度不适可能来自润滑不足或太快——\textbf{停下、补润滑、放慢}。
  \item \textbf{反复出血、持续疼痛、肛裂不愈}：就医（可查肛裂、痔疮、感染），别自行硬扛。
\end{itemize}

\section{误区}

\begin{enumerate}[leftmargin=2em,itemsep=0.3em]
  \item 「洗干净就安全」——干净挡不住"黏膜薄 + 无自润滑"，仍属高风险。
  \item 「不痛就安全 / 痛了忍忍就好」——疼痛是信号，微裂持续存在会持续抬升感染风险。
  \item 「肛交不伤身体」——反复微裂是真实伤害，需防护与必要时就医。
\end{enumerate}

\vspace{1em}
\noindent\textbf{一句话总结：} 肛交的安全公式 = \textbf{大量水/硅基润滑 + 安全套全程 + 前→后必换套 + 甲肝乙肝疫苗兜底（高危加 PrEP）}；把"无自润滑 + 黏膜薄"这两点记住，防护就不会漏。

\chapter{常见困惑深挖：很多"问题"其实是"正常"}
\label{chap:concerns}

很多让人尴尬的"情况"，其实处在\textbf{正常范围内}，或属于\textbf{可治疗、很常见}的问题——先别自我怀疑、也别羞耻。这一章把最常被问到的几类讲清楚。

\section{一、男性：勃起困难（ED）与早泄（PE）}

\subsection*{勃起功能障碍（ED）}
\begin{itemize}[leftmargin=2em,itemsep=0.3em]
  \item \textbf{原因常是两类叠加}：心理（焦虑、压力、关系紧张、表现焦虑）+ 生理（血管、神经、激素、某些药物）。
  \item \textbf{偶发很常见、属正常}；若\textbf{持续数周至数月且影响生活}，才值得就医。
  \item \textbf{可治疗}：治疗基础病（糖尿病/高血压/血脂）、在医生指导下使用 PDE5 抑制剂（如西地那非，\textbf{需处方、注意禁忌}）、心理与伴侣沟通。
  \item \textbf{红旗}：ED 有时是\textbf{心血管问题的早期信号}，医生常会顺带评估心血管风险。
\end{itemize}

\subsection*{早泄（PE）}
\begin{itemize}[leftmargin=2em,itemsep=0.3em]
  \item 分\textbf{原发}（一直如此）与\textbf{继发}（后来才出现，可能与其他因素有关）。
  \item \textbf{处理}：行为法（停--起始法、挤压法）、局部麻醉剂、心理（减压、与伴侣沟通）、在医生指导下使用处方药（部分 SSRI）。
  \item \textbf{关键}：别把它当"个人缺陷"——它是\textbf{常见、可管理}的，伴侣的理解与沟通本身就是治疗的一部分。
\end{itemize}

\section{二、女性：高潮困难与性疼痛}

\subsection*{高潮困难（anorgasmia）}
\begin{itemize}[leftmargin=2em,itemsep=0.3em]
  \item \textbf{多数女性需要直接阴蒂刺激}才到高潮（见\ref{chap:outer-pleasure}），"光插入不够"非常普遍，\textbf{不代表异常}。
  \item \textbf{常见原因}：焦虑/分心、关系因素、激素变化、药物（如部分抗抑郁药）、盆底紧张。
  \item \textbf{处理}：充分唤起与阴蒂刺激、与伴侣沟通偏好、必要时就医（评估激素/药物/盆底/心理）。
\end{itemize}

\subsection*{性疼痛（dyspareunia）/ 阴道痉挛（vaginismus）}
\begin{itemize}[leftmargin=2em,itemsep=0.3em]
  \item \textbf{疼痛原因多样}：感染、干涩、激素（围绝经期）、子宫内膜异位、盆底肌过度紧张、阴道痉挛。
  \item \textbf{处理}：润滑剂、局部激素（医生指导）、盆底治疗、心理、循序渐进的扩张训练、\textbf{找出根本原因}。
  \item \textbf{原则}：\textbf{疼不要忍}——反复疼痛应就医排查，而不是"习惯了就好"。
\end{itemize}

\section{三、双方：性欲差异（libido gap）}

\begin{itemize}[leftmargin=2em,itemsep=0.3em]
  \item \textbf{两人欲望不一样，很常见}，不等于"谁有问题"，更不等于谁不爱谁。
  \item 若\textbf{长期显著下降并造成困扰}，才考虑是否为\textbf{HSDD（女性性兴趣/欲望减退障碍）}，可就医评估（激素、药物、心理、关系）。
  \item \textbf{处理核心是沟通}（见\ref{chap:communication}）：谈感受、谈边界、找共同舒适的频率，\textbf{而非互相贬低或指责}。
\end{itemize}

\section{四、何时该去看医生（红旗信号）}

\begin{itemize}[leftmargin=2em,itemsep=0.3em]
  \item 反复/持续的\textbf{疼痛、出血、异常分泌物或体味}。
  \item \textbf{持续}的勃起困难、早泄并造成困扰。
  \item 性欲\textbf{长期大幅下降}并影响关系或情绪。
  \item 伴随\textbf{情绪困扰、创伤经历}——找正规心理/性医学科。
  \item 任何你"拿不准"的身体变化——\textbf{问医生比猜更靠谱}。
\end{itemize}

\section{一页要点}

\begin{tabularx}{\textwidth}{lCCX}
\toprule
情况 & 常见原因 & 可先做 & 何时就医 \\
\midrule
勃起困难 & 心理+血管/激素/药物 & 减压、沟通 & 持续数月、伴晨勃消失 \\
早泄 & 原发/继发、紧张 & 行为法、沟通 & 长期困扰 \\
高潮困难 & 需阴蒂刺激/心理/药物 & 前戏+阴蒂刺激 & 长期困扰、情绪受影响 \\
性疼痛 & 感染/干涩/激素/盆底 & 润滑、别忍 & 反复疼痛、出血 \\
性欲差异 & 关系/压力/激素/药物 & 沟通找原因 & 长期显著下降并困扰 \\
\bottomrule
\end{tabularx}

\vspace{1em}
\noindent\textbf{一句话总结：} 绝大多数"性方面的困扰"都落在\textbf{正常范围}或\textbf{常见、可治疗}的区间——\textbf{别羞耻、别硬扛、先沟通，持续或疼痛就找正规医生}，这是最省时间也最安全的路。

\chapter{术语表与全书总览}
\label{chap:glossary}

把全书出现的关键术语整理成中英对照，方便日后检索；再把整本书串一遍，方便你按主题快速翻回对应章节。

\section{核心术语中英对照}

\begin{tabularx}{\textwidth}{l C X}
\toprule
术语 & 英文 & 一句话 \\
\midrule
前戏 & foreplay & 唤起阶段的爱抚、亲吻、触摸 \\
边缘控制 & edging & 接近高潮时放慢/退回，再推上去 \\
同意 & consent & 持续、自愿、可撤回，不只是"没说" \\
偷摘安全套 & stealthing & 未经同意移除/不使用安全套 \\
长效可逆避孕 & LARC & 一次放置、多年有效的避孕（IUD/埋植） \\
宫内节育器 & IUD & 放入子宫的长效避孕装置 \\
性传播感染 & STI & 经性接触传播的感染总称 \\
疱疹 & HSV & 可口/生殖器互染，皮肤接触传播 \\
人乳头瘤病毒 & HPV & 与宫颈癌/口咽癌相关，疫苗可防 \\
牙科薄膜 & dental dam & 口交时的口腔屏障工具 \\
尿路感染 & UTI & 性行为后女性常见，事后排尿可降风险 \\
性疼痛 & dyspareunia & 性交疼痛，需找根因，别忍 \\
阴道痉挛 & vaginismus & 盆底肌不自主紧张致难以进入 \\
高潮困难 & anorgasmia & 长期难以达高潮，多与刺激/心理/药物相关 \\
勃起功能障碍 & ED & 持续勃起困难，常心理+生理，可治疗 \\
早泄 & PE & 射精过早泄且困扰，常见可管理 \\
性欲减退障碍 & HSDD & 性欲长期显著下降并困扰 \\
性欲 & libido & 对性的兴趣与欲望，因人而异 \\
暴露前预防 & PrEP & 高危者用药预防 HIV \\
暴露后预防 & PEP & 无保护暴露后 72h 内紧急用药 \\
紧急避孕 & emergency contraception & 事后 72h/120h 内的补救措施 \\
唤起伸展 & tenting & 唤起时阴道伸展增加空间 \\
性反应周期 & sexual response cycle & 兴奋/平台/高潮/消退四阶段 \\
\bottomrule
\end{tabularx}

\section{全书总览（按主题串一遍）}

\begin{enumerate}[leftmargin=2em,itemsep=0.3em]
  \item \textbf{为什么要读这本书}——定位与用法。
  \item \textbf{五个方面的准备}——认知、安全、沟通、健康、心态。
  \item \textbf{认知准备}——先懂 basics，破除谣言。
  \item \textbf{安全防护}——最实际的一章（安全套/润滑/检测）。
  \item \textbf{沟通与同意}——不可逾越的底线。
  \item \textbf{健康维护}——把风险挡在前面。
  \item \textbf{心态准备}——从容比技巧更重要。
  \item \textbf{去哪里获取可靠信息}——权威渠道清单。
  \item \textbf{深入篇：高频疑问解答}。
  \item \textbf{性传播感染详解}——各 STI 的循证说明。
  \item \textbf{关系与现实}——小白最该知道的几件事。
  \item \textbf{见面了怎么办}——从邀约到在场的清单。
  \item \textbf{男性视角的补充}——同样重要的一半。
  \item \textbf{到了私密场所后}——以酒店为例的具体清单。
  \item \textbf{床上的事}——亲密过程的健康框架。
  \item \textbf{性反应周期与注意事项}——四阶段的生理过程。
  \item \textbf{亲密行为的类型与健康防护}——四类行为总览。
  \item \textbf{非插入类行为详解}——完整、合理、不遗漏。
  \item \textbf{非插入类行为的愉悦技巧}——从"不危险"到"很愉悦"。
  \item \textbf{亲密关系沟通}——从"会说话"到"更亲密"。
  \item \textbf{避孕方法详解}——怎么选才靠谱。
  \item \textbf{口交、阴道性交、肛交：逐类深入}——三类对照。
  \item \textbf{阴道性交详解}——单独深入。
  \item \textbf{口交详解}——咽部/黏膜感染与屏障。
  \item \textbf{肛交详解}——无自润滑+黏膜薄的防护铁律。
  \item \textbf{常见困惑深挖}——很多"问题"其实是"正常"。
  \item \textbf{术语表与全书总览}——本条。
  \item \textbf{出发前的清单}——逐项确认的实操清单。
\end{enumerate}

\section{去哪里查权威信息}

\begin{itemize}[leftmargin=2em,itemsep=0.3em]
  \item 国际：\textbf{WHO（who.int）}、\textbf{CDC（cdc.gov）} 的性健康/STI 页面。
  \item 国内：\textbf{国家卫生健康委员会}、\textbf{中国疾控中心} 的官方发布。
  \item 就医：正规医院的\textbf{性医学科、男科/泌尿外科、妇科、感染科、心理科}。
  \item 原则：\textbf{以权威发布为准，别把网络偏方当结论}；拿不准就问医生。
\end{itemize}

\vspace{1em}
\noindent\textbf{全书一句话：} 性健康 = \textbf{知情 + 自愿 + 愉悦 + 安全}，四者缺一不可；这本书把"认知、防护、沟通、行为、方法、困惑"串成一条可执行的线，剩下的交给你的判断与医生的专业。

\chapter{附录：全书知识框架与补充索引}

前面各章是写给读者的内容，本附录换个视角：给出全书的\textbf{知识骨架与索引地图}，既方便按图索骥、也留作日后逐步补充的框架。每个节点标注了核心问题、关键概念、证据等级与覆盖状态。

\section{框架依据与图例}

\begin{tabularx}{\linewidth}{X X}
\toprule
\textbf{依据} & \textbf{说明} \\
\midrule
WHO 性健康定义（2002） & 与性有关的\textbf{身体、情感、心理、社会}四方面完好状态，不仅是无疾病；含"积极、尊重、愉悦、安全、免于胁迫歧视暴力" \\
生物-心理-社会模型 & 任何性问题都从\textbf{生物、心理、社会/关系}三层面看，避免单一归因 \\
性权利框架 & 性自由、平等、隐私、免于暴力、性健康权；对应同意、偷拍、影像传播等议题 \\
\bottomrule
\end{tabularx}

\vspace{0.5em}
\noindent\textbf{图例}：证据等级 \textbf{1} = 系统综述/权威指南（WHO/CDC/ACOG/EAU/NICE）；\textbf{2} = RCT/大型队列；\textbf{3} = 专家共识/权威教材；\textbf{4} = 科普共识/观察描述。覆盖状态 \textbf{已} = 已覆盖；\textbf{半} = 部分覆盖；\textbf{待} = 待补充。章号对应正文。

\section{A 基础认知}

\begin{tabularx}{\linewidth}{C X c c}
\toprule
\textbf{节点} & \textbf{核心问题 / 关键概念} & \textbf{证据} & \textbf{状态} \\
\midrule
A1 定义与权利 & 性健康不只是"没病"；性权利有哪些 & 1 & 已 \\
A2 女性解剖 & 内外生殖器；阴蒂/尿道/阴道位置；尿道短易 UTI & 3 & 已 \\
A3 男性解剖 & 阴茎/睾丸/附睾/前列腺；勃起射精生理 & 3 & 半 \\
A4 性反应周期 & 兴奋/平台/高潮/消退；唤起需时、阴蒂高潮 & 2 & 已 \\
A5 性取向与性别 & 取向与性别认同的区别；常见误解 & 3 & 待 \\
A6 性心理 & 性欲如何产生；心理唤起先于生理 & 3 & 半 \\
A7 生命周期总览 & 一生的性健康重点如何随年龄变化 & 3 & 已 \\
\bottomrule
\end{tabularx}

\section{B 健康与防护}

\begin{tabularx}{\linewidth}{C X c c}
\toprule
\textbf{节点} & \textbf{核心问题 / 关键概念} & \textbf{证据} & \textbf{状态} \\
\midrule
B1 STI 总论 & 病原谱、传播途径、无症状与筛查 & 1 & 已 \\
B2 重点 STI & HIV/HPV/疱疹/梅毒/淋病/衣原体/乙肝 & 1 & 已 \\
B3 避孕方法 & 完美 vs 典型失败率；按需求选 & 1 & 已 \\
B4 紧急避孕 & 72h/120h/铜 IUD 5 天 & 1 & 已 \\
B5 屏障工具 & 安全套用法、口用套/薄膜、润滑选择 & 1 & 已 \\
B6 暴露前后预防 & PrEP/PEP 适用与就医 & 1 & 半 \\
B7 疫苗 & HPV/乙肝/甲肝（肛交者） & 1 & 半 \\
B8 女性非 STI 问题 & BV/念珠菌/UTI 识别与就医 & 2 & 半 \\
B9 男性常见问题 & ED/早泄；处方药与就医指征 & 1 & 半 \\
B10 检测与筛查流程 & 查什么、窗口期、去哪查、阳性怎么办 & 1 & 待 \\
\bottomrule
\end{tabularx}

\section{C 行为与实践}

\begin{tabularx}{\linewidth}{C X c c}
\toprule
\textbf{节点} & \textbf{核心问题 / 关键概念} & \textbf{证据} & \textbf{状态} \\
\midrule
C1 行为分类 & 插入/非插入；体液 vs 皮肤接触风险 & 1 & 已 \\
C2 非插入类 & 阴蒂为核心；慢+润滑、边缘控制 & 2 & 已 \\
C3 阴道性交 & 唤起、润滑、阴蒂叠加、防 UTI & 2 & 已 \\
C4 口交 & 咽部感染、薄膜/口用套、HSV-1 & 1 & 已 \\
C5 肛交 & 无自润滑、黏膜薄、换套、水/硅基润滑 & 1 & 已 \\
C6 愉悦元技巧 & 慢、润滑、读反应、开口问、高潮非 KPI & 2 & 已 \\
C7 姿势与场景 & 按需求选；孕期/初次调整 & 4 & 已 \\
C8 常见困惑 & 高潮困难/性疼痛/性欲差异；红旗信号 & 2 & 已 \\
\bottomrule
\end{tabularx}

\section{D 关系、同意与伦理}

\begin{tabularx}{\linewidth}{C X c c}
\toprule
\textbf{节点} & \textbf{核心问题 / 关键概念} & \textbf{证据} & \textbf{状态} \\
\midrule
D1 同意 & 有效同意要素、可撤回、stealthing & 1 & 已 \\
D2 沟通 & NVC 四步、时机、倾听 & 3 & 已 \\
D3 约会与现实 & 邀约到在场的清单、防偷拍 & 4 & 已 \\
D4 关系形态 & 长期/异地/开放关系的边界协商 & 3 & 待 \\
D5 法律伦理 & 隐私、偷拍、影像传播的法规边界 & 1 & 半 \\
D6 数字时代的同意 & 私密影像、AI 合成、网络骚扰 & 3 & 待 \\
\bottomrule
\end{tabularx}

\section{E 生命周期与特殊人群}

\begin{tabularx}{\linewidth}{C X c c}
\toprule
\textbf{节点} & \textbf{核心问题 / 关键概念} & \textbf{证据} & \textbf{状态} \\
\midrule
E1 青春期 & 身体变化、保护习惯、可靠信息 & 1 & 已 \\
E2 孕期 & 安全性、红线、姿势、干涩、引产 & 1 & 已 \\
E3 产后 & 恢复时机、避孕陷阱、LAM 三条件 & 1 & 已 \\
E4 更年期 & GSM 干涩、局部雌激素、避孕时限 & 1 & 已 \\
E5 老年 & 需求、ED、药物、防护 & 2 & 已 \\
E6 残障/慢性病 & 身体限制下的亲密与用药 & 3 & 待 \\
E7 LGBTQ+ 健康 & 针对性筛查、去污名 & 2 & 待 \\
\bottomrule
\end{tabularx}

\section{F 求助、资源与工具}

\begin{tabularx}{\linewidth}{C X c c}
\toprule
\textbf{节点} & \textbf{核心问题 / 关键概念} & \textbf{证据} & \textbf{状态} \\
\midrule
F1 就医路径 & 症状对应科室、如何开口 & 3 & 待 \\
F2 红旗信号 & 剧痛、大出血、发热、神经症状须就医 & 1 & 半 \\
F3 权威信息源 & WHO/CDC/ACOG/NHS/中国疾控/卫健委 & 1 & 已 \\
F4 术语表 & 中英术语对照 & --- & 已 \\
F5 检测与疫苗地图 & 本地化检测/疫苗渠道 & 3 & 待 \\
F6 更新维护 & 每次补充记录日期/节点/来源 & --- & 待 \\
\bottomrule
\end{tabularx}

\section{怎么补充才科学}

\begin{enumerate}[leftmargin=2em,itemsep=0.3em]
  \item \textbf{一个节点一段}：先写结论，再写"为什么"（机制/证据），最后写"怎么做"。
  \item \textbf{标注证据等级}：新结论注明来源机构或文献；无权威来源者标"观点/共识"，不写成定论。
  \item \textbf{数字要核对}：失败率、时限、剂量等关键数值，写入前联网核对权威来源（同正文做法）。
  \item \textbf{中国语境}：涉及法律、就医、疫苗可及性时以中国法规与本地服务为准，注明地区差异。
  \item \textbf{避免重复}：补充前先看状态列，已覆盖的节点按需"深化"而非新建。
  \item \textbf{配图}：流程/分层/曲线类节点按全书 TikZ 风格补图，保持风格统一。
\end{enumerate}

\chapter{出发前的清单}

把前面的内容浓缩为一份可操作的清单。逐项确认，就能大幅降低风险。

\begin{enumerate}[leftmargin=2em,itemsep=0.4em]
  \item 我了解基本的生理与避孕常识，不靠谣言做决定。
  \item 我准备了合格、在有效期内、尺寸合适的安全套，并会用。
  \item 我备好水基润滑剂，知道油基会破坏乳胶。
  \item 我和伴侣已经谈过避孕、检测和各自的边界。
  \item 我确认双方都清醒、自愿，且可随时喊停。
  \item 我了解 HPV 疫苗，并考虑接种与定期筛查。
  \item 我知道无保护行为后的紧急避孕与就医途径。
  \item 我知道去哪里查权威信息，也知道遇到问题该找医生。
\end{enumerate}

\chapter*{结语}
\addcontentsline{toc}{chapter}{结语}

性健康不是一门靠"顿悟"的学问，而是可以系统学习、理性准备的知识。
愿这本小册子帮你建立起安全、自愿、知性的观念：既保护自己，也尊重他人。
当你需要更深入的专业内容时，请翻阅同书库的《性医学基础》，
或直接向正规医疗机构的专业人员求助。

\vspace{2em}

\noindent\textbf{主要资料来源：}
\begin{itemize}[leftmargin=2em]
  \item World Health Organization (WHO). \emph{Family planning/contraception methods}; \emph{Medical eligibility criteria for contraceptive use}, 6th ed.
  \item U.S. Centers for Disease Control and Prevention (CDC). \emph{Primary Prevention Methods — STI Treatment Guidelines}; \emph{Condom Use: An Overview}; \emph{The Lowdown on How to Prevent STDs}.
  \item 国家卫生健康委员会及中国疾病预防控制中心关于 HPV 疫苗接种的相关公告（以最新发布为准）。
  \item European Association of Urology (EAU). \emph{Guidelines on Male Sexual Dysfunction}; ICD-11"早期射精"定义（第14章）。
  \item 科技日报、中国新闻网：酒店针孔摄像头识别与防范专家建议（藏在镜子背面、插座、空调出风口、烟雾报警器等处；手机"红点法"局限；手电查反光、查 WiFi 陌生设备、物理遮挡）（第15章）。
  \item 全国人大《中华人民共和国民法典》第 1033 条、《治安管理处罚法》、《刑法》第 283 条关于隐私空间与窃照器材的规定（第15章）。
  \item Non-consensual condom removal ("stealthing"): PMC\,6306234（墨尔本性健康诊所研究）；Drouin 等，\emph{Psychology of Violence} (2024)（第15章）。
  \item CDC. \emph{Clinical Overview of Urinary Tract Infections}（性行为引入细菌、女性尿道较短风险更高）；\emph{Condom Use: An Overview}（第16章）。
  \item WHO. \emph{Sexual health}（2006 工作定义：积极、尊重、愉悦且安全、免于胁迫）；\emph{Redefining sexual health for benefits throughout life} (2022)（第16章）。
  \item W.\,H. Masters \& V.\,E. Johnson. \emph{Human Sexual Response} (1966)：兴奋期/平台期/高潮期/消退期四阶段模型（第17章）。
  \item WHO. \emph{Comprehensive sexuality education}（科学准确、适龄的身体与关系知识，帮助保护健康、尊重他人）（第17章）。
  \item CDC. \emph{How to Prevent STIs}; \emph{About STI Risk and Oral Sex}; \emph{How to Use a Dental Dam}; \emph{Condom Use: An Overview}（屏障工具、换套、口交膜与杀精剂禁忌）（第18章）。
  \item CDC. \emph{MMWR Recommendations—Opportunistic Pathogens}（避免口--粪暴露、换套、手部清洁）；WHO. \emph{Condoms}（直肠无自然润滑、肛交需润滑）（第18章）。
  \item CDC. \emph{Appendix A: Known or Possible HIV Exposure Scenarios}（接吻、相互手刺激属可忽略风险，nPEP 不常规推荐，除非可见血 + 黏膜破损）；\emph{How HIV Spreads}（HIV 仅经特定体液经黏膜/破损皮肤传播）（第19章）。
  \item CDC. \emph{Preventing HIV with Condoms}（安全套用于共用器具高效防 HIV 及部分 STI）；\emph{How to Prevent STIs}（皮肤接触传播）（第19章）。
  \item 女性高潮与阴蒂刺激依赖：大样本与综述研究（多数女性需直接阴蒂刺激达高潮；仅阴道插入缺少阴蒂刺激时高潮率明显更低）——综合 Masters \& Johnson 性反应模型与后续大样本调查（第19章）。
  \item CDC. \emph{How to Use a Dental Dam}; \emph{About STI Risk and Oral Sex}（口--口/口--生殖器黏膜传播、口腔屏障工具、有活动性溃疡/疱疹时避免）（第19章）。
  \item M.\,B.\,Rosenberg. \emph{Nonviolent Communication: A Language of Life}（观察--感受--需要--请求四步框架，请求与指责的区分）（第20章）。
  \item WHO. \emph{Sexual health}（积极、尊重、愉悦且安全的性体验，免于胁迫与暴力）（第20章）。
  \item American College of Obstetricians and Gynecologists (ACOG); J.\,Trussell 等. \emph{Typical-Use and Perfect-Use Contraceptive Failure Rates}（各方法典型/完美使用失败率分层）（第21章）。
  \item WHO. \emph{Medical Eligibility Criteria for Contraceptive Use}, 6th ed.（口服药适应证/禁忌、避孕方法选择）（第21章）。
  \item CDC. \emph{Emergency Contraception}（左炔诺孕酮 72h 内最佳、乌利司坦 120h、铜 IUD 事后 5 天内最可靠）（第21章）。
  \item CDC. \emph{Pharyngeal Gonorrhea/Chlamydia}; \emph{About Hepatitis A and B}（有肛交者建议接种甲肝/乙肝疫苗）（第22章）。
  \item WHO. \emph{Condoms}（直肠无自然润滑、肛交需充分水基/硅基润滑；油基破坏乳胶）（第22章）。
  \item CDC. \emph{Clinical Overview of Urinary Tract Infections}（女性尿道短、性行为引入细菌、事后排尿/多饮水降低 UTI 风险）；\emph{Vaginal Douching}（阴道冲洗破坏菌群，不推荐）（第22章）。
  \item W.\,H.\,Masters \& V.\,E.\,Johnson. \emph{Human Sexual Response} (1966)（唤起期"tenting"伸展、高潮非必然）；ACOG 关于 \emph{Vaginismus / dyspareunia} 的临床说明（剧痛就医指征）（第23章）。
  \item WHO. \emph{Condoms}; CDC. \emph{Lubricants and Condoms}（水基/硅基可配乳胶套、油基破坏乳胶）（第23章）。
  \item WHO. \emph{Human papillomavirus (HPV)}（HPV 与宫颈癌/口咽癌关联、疫苗可防；HIV 经口传播极低的判断与\ref{chap:sti}一致）（第24章）。
  \item CDC. \emph{PrEP (Pre-Exposure Prophylaxis) for HIV Prevention}（高危者 PrEP 显著降低 HIV 感染风险）（第25章）。
  \item EAU. \emph{Guidelines on Erectile Dysfunction / Premature Ejaculation}（ED/PE 循证处理、PDE5 抑制剂需处方、行为法）；NICE / AUA 关于 \emph{dyspareunia、vaginismus} 与 HSDD 的说明（第26章）。
  \item NHS. \emph{Is it safe to have sex in pregnancy?}；Tommy's \emph{Sex in pregnancy}（孕期性生活一般安全、需避免的情况：前置/低位胎盘·破水·多次流产·不明原因出血；推荐并排侧卧、水基润滑；口唇疱疹避免口--生殖器；无证据表明性交会引产、假宫缩非临产；28 周后避免长时间仰卧）（第25章）。
  \item 阴道瓣（hymen）解剖与初次性行为：阴道口黏膜褶皱、中间有孔、形态多样；初次不一定出血或疼痛；疼痛/出血多源于紧张·润滑不足·动作过快（科普共识）（第25章）。
  \item NHS. \emph{Vaginal oestrogen / HRT}（更年期局部雌激素低剂量、不增加乳腺癌风险、无需配孕激素、可按需长期使用、约 3 个月起效；睾酮仅专科特定情况；停经后避孕时限 50 岁以下 2 年/以上 1 年、55 岁后极罕见）（第26章）。
  \item WHO. \emph{Family planning/contraception methods}（安全套唯一既避孕又防 STI、现代避孕不致不孕、青少年可安全使用）；\emph{Medical eligibility criteria / Oral contraceptives}（产后避孕时机：纯母乳 6 个月/部分母乳 6 周、迷你 pill 产后即可）。LAM 哺乳闭经避孕法三条件（纯母乳·月经未恢复·6 个月内）（第26章）。
\end{itemize}

\begin{center}
\itshape 本文为健康教育科普，不构成医疗建议；具体诊疗请遵医嘱。
\end{center}

\end{document}
